\section{Differential criteria for smoothness and regularity}
\label{section:0.22}

\oldpage[0\textsubscript{IV-1}]{182}
The main results of this section are:
\begin{enumerate}
  \item[a)] The criteria \sref{0.22.2.2} and \sref{0.22.5.8}, which give necessary and sufficient conditions for a complete Noetherian local ring $A$ containing a field $k$ to be formally smooth over $k$; the criterion \sref{0.22.5.8} is stated without involving differential notions, and completes \sref{0.19.6.6}; the statement \sref{0.22.2.2}, of a slightly more technical nature, permits giving a \emph{classification}, for fixed $k$, of the local rings $A$ in question \sref{0.22.2.5}: they are determined by their dimension $n$ and their residue field $\mathrm{K}$, under the sole condition $n\geq\rg_\mathrm{K}\Upsilon_{\mathrm{K}/k}$.
  \item[b)] The theorem \sref{0.22.3.3}, and the corollary \sref{0.22.7.6} of the Jacobian criterion of Nagata \sref{0.22.7.3}, which give conditions allowing one to assert that a \emph{localized} ring of a complete Noetherian local ring $A$ containing a field $k$ is geometrically regular over $k$; they will play an important role in the study of the ``fine'' properties of local rings carried out in chap.~IV, \textsection~7. We do not possess for the moment any proof of \sref{0.22.7.6} independent of the Jacobian criterion of Nagata.
  \item[c)] The Jacobian criterion of Zariski \sref{0.22.6.7} and its variants, which are easy consequences of \sref{0.20.7.8}.
\end{enumerate}

\subsection{Lifting of formal smoothness}
\label{subsection:0.22.1}

\oldpage[0\textsubscript{IV-1}]{183}
\begin{theorem}[22.1.1]
\label{0.22.1.1}
Let $s:\Lambda\to\mathrm{A}$, $u:\mathrm{A}\to\mathrm{B}$ be two ring homomorphisms, $\mathfrak{J}$ an ideal of $\mathrm{B}$, $\mathrm{C}$ the quotient ring $\mathrm{B}/\mathfrak{J}$.
Suppose that $\mathrm{C}$ is a $\Lambda$-algebra formally smooth for the discrete topology.
\begin{enumerate}
  \item[\rm{(i)}] \emph{Suppose the following conditions are satisfied:}
  \begin{enumerate}
    \item[$\alpha$)] $\mathfrak{J}/\mathfrak{J}^2$ \emph{is a projective $\mathrm{C}$-module.}
    \item[$\beta$)] \emph{The canonical homomorphism $\mathbf{S}^*_\mathrm{C}(\mathfrak{J}/\mathfrak{J}^2)\to\gr^*_\mathfrak{J}(\mathrm{B})$ \sref{0.19.5.2} is bijective.}
    \item[$\gamma$)] \emph{The characteristic homomorphism $\chi:\Upsilon_{\mathrm{C}/\Lambda/\Lambda}\to\mathfrak{J}/\mathfrak{J}^2$ \sref{0.20.6.20} is injective.}
    \item[$\delta$)] \emph{The $\mathrm{C}$-module $\Omega^1_{\mathrm{B}/\Lambda}\otimes_\mathrm{B}\mathrm{C}$ is projective.}
  \end{enumerate}
  \emph{Then $\mathrm{B}$, equipped with the $\mathfrak{J}$-preadic topology, is a formally smooth $\Lambda$-algebra.}
  \item[\rm{(ii)}] \emph{Conversely, suppose that $\mathrm{B}$ is a $\Lambda$-algebra formally smooth for the $\mathfrak{J}$-preadic topology, and that $\mathrm{A}$ is a $\Lambda$-algebra formally smooth for the discrete topology.
  Then the conditions $\alpha$), $\beta$), $\gamma$), $\delta$) of \rm{(i)} are satisfied.}
\end{enumerate}
\end{theorem}

\begin{proof}
By virtue of the exact sequence \sref{0.20.6.22.1} (which is applicable, since $\mathrm{B}/\mathfrak{J}^2$ is a $\Lambda$-trivial extension of $\mathrm{C}$ by $\mathfrak{J}/\mathfrak{J}^2$, $\mathrm{C}$ being assumed to be a formally smooth $\Lambda$-algebra \sref{0.19.4.4}), the condition $\gamma$) is equivalent to the following:
\begin{enumerate}
  \item[$\gamma'$)] \emph{The homomorphism $u_{\mathrm{B}/\Lambda/\Lambda}\otimes 1_\mathrm{C}:\Omega^1_{\Lambda/\Lambda}\otimes_\mathrm{A}\mathrm{C}\to\Omega^1_{\mathrm{B}/\Lambda}\otimes_\mathrm{B}\mathrm{C}$ is injective.}
\end{enumerate}

(i) The conditions $\alpha$) and $\beta$) imply that $\mathrm{B}$ is a $\Lambda$-algebra formally smooth for the $\mathfrak{J}$-preadic topology.
It amounts to the same thing to say that $\mathrm{B}$ is an $\mathrm{A}$-algebra formally smooth \emph{relative} to $\Lambda$, or an $\mathrm{A}$-algebra formally smooth relative to $\Lambda$.
To see that $\mathrm{B}$ is an $\mathrm{A}$-algebra formally smooth for the $\mathfrak{J}$-preadic topology, it suffices therefore, by virtue of \sref{0.20.7.2}, to prove that the continuous homomorphism of topological $\mathrm{B}$-modules $u_{\mathrm{B}/\Lambda/\Lambda}:\Omega^1_{\Lambda/\Lambda}\otimes_\mathrm{A}\mathrm{B}\to\Omega^1_{\mathrm{B}/\Lambda}$ is \emph{formally invertible on the left}.
But since $\mathrm{A}$ is a $\Lambda$-algebra formally smooth, the topological $\mathrm{B}$-module $\Omega^1_{\mathrm{B}/\Lambda}$ is formally projective \sref{0.20.4.9}, and its topology is deduced from that of $\mathrm{B}$ \sref{0.20.4.5}.
By virtue of \sref{0.19.1.9}, it suffices therefore, for $u_{\mathrm{B}/\Lambda/\Lambda}$ to be formally invertible on the left, that $u_{\mathrm{B}/\Lambda/\Lambda}\otimes 1_\mathrm{C}$ be invertible on the left.
But by virtue of $\gamma'$), this last map is injective, and on the other hand
\[
  0\to\Omega^1_{\Lambda/\Lambda}\otimes_\mathrm{A}\mathrm{C}\to\Omega^1_{\mathrm{B}/\Lambda}\otimes_\mathrm{B}\mathrm{C}\to\Omega^1_{\mathrm{B}/\Lambda}\otimes_\mathrm{B}\mathrm{C}\to 0.
\]

Finally, by virtue of $\delta$), the $\mathrm{C}$-module $\Omega^1_{\mathrm{B}/\Lambda}\otimes_\mathrm{B}\mathrm{C}$ is projective, so the preceding exact sequence is split, which finishes proving (i).

\oldpage[0\textsubscript{IV-1}]{184}
(ii) If $\mathrm{B}$ is a $\Lambda$-algebra formally smooth for the $\mathfrak{J}$-preadic topology, $\mathrm{B}$ is an $\mathrm{A}$-algebra formally smooth for the $\mathfrak{J}$-preadic topology as well, since $\mathrm{A}$ is a $\Lambda$-algebra formally smooth for the discrete topology; the conditions $\alpha$) and $\beta$) follow from \sref{0.19.5.4}.
It amounts to the same thing to say that $\mathrm{B}$ is an $\mathrm{A}$-algebra formally smooth relatively, or an $\mathrm{A}$-algebra formally smooth relative to $\Lambda$.
To see that $\mathrm{B}$ is a $\Lambda$-algebra formally smooth, it suffices therefore, by virtue of \sref{0.20.7.2}, to prove that the continuous homomorphism of topological $\mathrm{B}$-modules $u_{\mathrm{B}/\Lambda/\Lambda}:\Omega^1_{\Lambda/\Lambda}\otimes_\mathrm{A}\mathrm{B}\to\Omega^1_{\mathrm{B}/\Lambda}$ is \emph{formally invertible on the left} \sref{0.19.1.7}, and \emph{a fortiori} injective.
Finally, $\Omega^1_{\mathrm{B}/\Lambda}$ is a $\mathrm{B}$-module formally projective \sref{0.20.4.9}, and consequently $\Omega^1_{\mathrm{B}/\Lambda}\otimes_\mathrm{B}\mathrm{C}$ is a projective $\mathrm{C}$-module \sref{0.19.2.4}; the equivalence of $c$) follows then from \sref{0.19.1.9}.
\end{proof}

\begin{corollary}[22.1.2]
\label{0.22.1.2}
Let $s:\Lambda\to\mathrm{A}$, $u:\mathrm{A}\to\mathrm{B}$ be two ring homomorphisms, $\mathfrak{m}$ a maximal ideal of $\mathrm{B}$, $\mathrm{K}=\mathrm{B}/\mathfrak{m}$ the quotient field.
Suppose that the $\Lambda$-algebras $\mathrm{A}$ and $\mathrm{K}$ are formally smooth for the discrete topology.
Then the following conditions are equivalent:
\begin{enumerate}
  \item[a)] $\mathrm{B}$ is an $\mathrm{A}$-algebra formally smooth for the $\mathfrak{m}$-preadic topology.
  \item[b)] \emph{The canonical homomorphism}
\end{enumerate}
\[
\label{0.22.1.2.1}
  \mathbf{S}^*_\mathrm{K}(\mathfrak{m}/\mathfrak{m}^2)\to\gr^*_\mathfrak{m}(\mathrm{B})
  \tag{22.1.2.1}
\]
\emph{is bijective, and the characteristic homomorphism}
\[
\label{0.22.1.2.2}
  \chi:\Upsilon_{\mathrm{K}/\Lambda/\Lambda}\to\mathfrak{m}/\mathfrak{m}^2
  \tag{22.1.2.2}
\]
\emph{is injective.}
\end{corollary}

\begin{proof}
This follows immediately from \sref{0.22.1.1}, the conditions $\alpha$) and $\delta$) being trivially satisfied here, since $\mathrm{K}$ is a field.
\end{proof}

\begin{remark}[22.1.3]
\label{0.22.1.3}
Suppose the hypotheses of \sref{0.22.1.2} are satisfied, and suppose in addition that $\mathfrak{m}/\mathfrak{m}^2$ is a $\mathrm{K}$-vector space of \emph{finite} rank.
Then, for $\mathrm{B}$ to be an $\mathrm{A}$-algebra formally smooth for the $\mathfrak{m}$-preadic topology, it suffices that the following conditions be satisfied:

\begin{env}[22.1.3.1]
\label{0.22.1.3.1}
Given a polynomial algebra $\mathrm{E}=\mathrm{K}[\mathrm{T}_1,\ldots,\mathrm{T}_n]$ with indeterminates, and an ideal $\mathfrak{n}$ de $\mathrm{E}$ generated by the $\mathrm{T}_i$ $(1\leq i\leq n)$, and the maximal ideal $\mathfrak{n}$ of $\mathrm{E}$, every $\Lambda$-homomorphism $v:\mathrm{B}\to\mathrm{E}/\mathfrak{b}$, which, by passing to quotients, gives the identity $\mathrm{K}\to\mathrm{K}$, factors as $\mathrm{B}\xrightarrow{w}\mathrm{E}/\mathfrak{n}^k\to\mathrm{E}/\mathfrak{b}$, where $w$ is a $\Lambda$-homomorphism.
\end{env}

\begin{env}[22.1.3.2]
\label{0.22.1.3.2}
If $\mathrm{F}$ is the trivial extension $\mathrm{D}_\mathrm{K}(\mathrm{K})$ (``algebra of dual numbers over $\mathrm{K}$''), $\rho:\mathrm{A}\to\mathrm{F}$ a ring homomorphism such that the composition $\mathrm{A}\xrightarrow{\rho}\mathrm{F}\to\mathrm{K}$ is the canonical homomorphism, then every $\Lambda$-homomorphism $v:\mathrm{B}\to\mathrm{K}$ which, by passing to quotients, gives the identity $\mathrm{K}\to\mathrm{K}$, factors as $\mathrm{B}\xrightarrow{w}\mathrm{F}\to\mathrm{K}$, where $w$ is a $\Lambda$-homomorphism (for the $\mathrm{A}$-algebra structure on $\mathrm{F}$ defined by $\rho$).
\end{env}

Indeed, it was seen \sref{0.19.5.5} that the condition \sref{0.22.1.3.1} implies that the canonical homomorphism \sref{0.22.1.2.1} is bijective.
By virtue of \sref{0.22.1.2}, it suffices then to see that $u_{\mathrm{B}/\Lambda/\Lambda}\otimes 1_\mathrm{K}:\Omega^1_{\Lambda/\Lambda}\otimes_\mathrm{A}\mathrm{K}\to\Omega^1_{\mathrm{B}/\Lambda}\otimes_\mathrm{B}\mathrm{K}$ is injective, or equivalently (since these are vector spaces over $\mathrm{K}$) that for every $\mathrm{K}$-vector space $\mathrm{L}$, the canonical homomorphism $\Der_\Lambda(\mathrm{B},\mathrm{L})\to\Der_\Lambda(\mathrm{A},\mathrm{L})$ \emph{is surjective} (by virtue of \sref{0.20.4.8}); it is clear that for this, it suffices that the canonical homomorphism $\Der_\Lambda(\mathrm{B},\mathrm{K})\to\Der_\Lambda(\mathrm{A},\mathrm{K})$ is surjective, and our assertion follows, by virtue of \sref{0.20.1.5}.

Note that the two conditions \sref{0.22.1.3.1} and \sref{0.22.1.3.2} are implied by the following:

\begin{env}[22.1.3.3]
\label{0.22.1.3.3}
For every local Artinian $\mathrm{A}$-algebra $\mathrm{E}$ with residue field equal to $\mathrm{K}$, and every nilpotent ideal $\mathfrak{r}$ of $\mathrm{E}$, every $\Lambda$-homomorphism $\mathrm{B}\to\mathrm{E}/\mathfrak{r}$ which, by passing to quotients, gives the identity $\mathrm{K}\to\mathrm{K}$, factors as $\mathrm{B}\xrightarrow{u}\mathrm{E}\to\mathrm{E}/\mathfrak{r}$, where $u$ is a $\Lambda$-homomorphism.
\end{env}
\end{remark}

\oldpage[0\textsubscript{IV-1}]{185}
\begin{proposition}[22.1.4]
\label{0.22.1.4}
Let $\varphi:\mathrm{A}\to\mathrm{B}$ be a local homomorphism of local Noetherian rings, $k$ (resp.~$\mathrm{K}$) the residue field of $\mathrm{A}$ (resp.~$\mathrm{B}$).
For $\mathrm{B}$ to be a formally smooth $\mathrm{A}$-algebra (for the preadic topologies), it is necessary and sufficient that for every local Artinian ring $\mathrm{C}$, with residue field isomorphic to $\mathrm{K}$, every local homomorphism $\mathrm{A}\to\mathrm{C}$ (making $\mathrm{C}$ a topological $\mathrm{A}$-algebra), and every nilpotent ideal $\mathfrak{J}$ of $\mathrm{C}$, every local $\mathrm{A}$-homomorphism $\mathrm{B}\to\mathrm{C}/\mathfrak{J}$ such that the homomorphism $\mathrm{K}\to\mathrm{K}$ obtained by passing to quotients is the identity, factors as $\mathrm{B}\xrightarrow{f}\mathrm{C}\to\mathrm{C}/\mathfrak{J}$, where $f$ is a local $\mathrm{A}$-homomorphism.
\end{proposition}

\begin{proof}
The condition being necessary by definition \sref{0.19.3.1}, let us show that it is sufficient.
Since the $\widehat{\mathrm{A}}$-local homomorphisms of $\widehat{\mathrm{B}}$ in this discrete $\widehat{\mathrm{A}}$-algebra come from, by extension, the $\mathrm{A}$-local homomorphisms of $\mathrm{B}$ in this algebra, we may restrict ourselves to the case where $\mathrm{A}$ and $\mathrm{B}$ are \emph{complete}, by virtue of \sref{0.19.3.6}.
When $\mathrm{A}$ is a \emph{field} $k$, the proposition follows from \sref{0.22.1.3.3}, where one takes $\Lambda$ to be the prime subfield of $k$, and \sref{0.19.6.1}.
In the general case, set $\mathrm{B}_0=\mathrm{B}\otimes_\mathrm{A}k=\mathrm{B}/\mathfrak{m}\mathrm{B}$ ($\mathfrak{m}$ the maximal ideal of $\mathrm{A}$), which is complete.
If $\mathrm{C}_0$ is a local Artinian $k$-algebra with residue field isomorphic to $\mathrm{K}$, $\mathfrak{J}_0$ a nilpotent ideal of $\mathrm{C}_0$, every $k$-local homomorphism $g_0:\mathrm{B}_0\to\mathrm{C}_0/\mathfrak{J}_0$ giving the identity on $\mathrm{K}$ by passing to quotients, gives by composition an $\mathrm{A}$-local homomorphism
\[
  g:\mathrm{B}\to\mathrm{B}_0\xrightarrow{g_0}\mathrm{C}_0/\mathfrak{J}_0
\]
which by hypothesis factors as $\mathrm{B}\xrightarrow{f}\mathrm{C}_0\to\mathrm{C}_0/\mathfrak{J}_0$, where $f$ is a local $\mathrm{A}$-homomorphism; but as $\mathrm{C}_0$ is a $k$-algebra, $f$ factors as $\mathrm{B}\to\mathrm{B}_0\xrightarrow{f_0}\mathrm{C}_0$, where $f_0$ is a local $k$-homomorphism.
We thus see that the hypotheses of the statement are still fulfilled when one substitutes $\mathrm{B}_0$ and $k$ for $\mathrm{B}$ and $\mathrm{A}$ respectively; therefore $\mathrm{B}_0$ is a $k$-algebra \emph{formally smooth} from what we have just seen.
Applying \sref{0.19.7.2} and \sref{0.19.7.1}, we see that there exists a local Noetherian complete ring $\mathrm{B}'$, a local homomorphism $\mathrm{A}\to\mathrm{B}'$ making $\mathrm{B}'$ a flat $\mathrm{A}$-module and an $\mathrm{A}$-algebra \emph{formally smooth}, and a $k$-isomorphism $u_0':\mathrm{B}_0'=\mathrm{B}'\otimes_\mathrm{A}k\xrightarrow{\sim}\mathrm{B}_0$.
Let $v_0'$ be the inverse isomorphism of $u_0'$; we now propose to show that there exists a local $\mathrm{A}$-homomorphism $v:\mathrm{B}\to\mathrm{B}'$ such that $v_0'$ comes from $v$ by passing to quotients.
For this, let us denote by $\mathfrak{n}$ and $\mathfrak{n}'$ the maximal ideals of $\mathrm{B}$ and $\mathrm{B}'$ respectively, and for every $j$, let $w_j:\mathrm{B}/(\mathfrak{n}^{j+1}+\mathfrak{m}\mathrm{B})\to\mathrm{B}'/(\mathfrak{n}'^{j+1}+\mathfrak{m}\mathrm{B}')$ be the $\mathrm{A}$-local homomorphism deduced from $v_0'$ by passing to quotients; it will suffice to form for each $j$ an $\mathrm{A}$-local homomorphism $v_j:\mathrm{B}/\mathfrak{n}^{j+1}\to\mathrm{B}'/\mathfrak{n}'^{j+1}$ such that the diagrams
\[
  \begin{tikzcd}
    \mathrm{B}/\mathfrak{n}^{j+1} \ar[r,"v_j"] & \mathrm{B}'/\mathfrak{n}'^{j+1} & \mathrm{B}/\mathfrak{n}^{j+1} \ar[r,"v_j"] & \mathrm{B}'/\mathfrak{n}'^{j+1}\\
    \mathrm{B}/\mathfrak{n}^j \ar[u] \ar[r,"v_{j-1}"] & \mathrm{B}'/\mathfrak{n}'^j \ar[u] & \mathrm{B}/(\mathfrak{n}^{j+1}+\mathfrak{m}\mathrm{B}) \ar[u] \ar[r,"w_j"] & \mathrm{B}'/(\mathfrak{n}'^{j+1}+\mathfrak{m}\mathrm{B}') \ar[u]
  \end{tikzcd}
\]

\oldpage[0\textsubscript{IV-1}]{186}
are commutative; $\mathrm{B}$ and $\mathrm{B}'$ being complete, the homomorphism $v=\varprojlim v_j$ will give the answer.
But the formation by induction of $v_j$ follows from the hypothesis on $\mathrm{B}$ and from the fact that $\mathfrak{n}'^j+(\mathfrak{n}'^{j+1}+\mathfrak{m}\mathrm{B}')=\mathfrak{n}'^j+\mathfrak{m}\mathrm{B}'$.
It follows then from \sref{0.19.7.1.4} that $v$ is an $\mathrm{A}$-\emph{isomorphism}, for $\mathrm{B}'$ is a flat $\mathrm{A}$-module, and $\mathrm{B}$ is complete for the $\mathfrak{m}$-preadic topology, the ideals $\mathfrak{m}^i\mathrm{B}$ being closed in $\mathrm{B}$ for the $\mathfrak{n}$-adic topology $(\mathbf{0}_\mathrm{I},~7.3.5)$.
\hfill Q.E.D.
\end{proof}

\subsection{Differential characterization of complete local algebras formally smooth over a field}
\label{subsection:0.22.2}

\begin{env}[22.2.1]
\label{0.22.2.1}
Let $k$ be a field, $\mathrm{P}$ its prime subfield, $\mathrm{A}$ a $k$-algebra that is a \emph{local} ring, $\mathfrak{m}$ its maximal ideal, $\mathrm{K}=\mathrm{A}/\mathfrak{m}$ its residue field.
As $\mathrm{K}$ is separable over $\mathrm{P}$, $\mathrm{K}$ is a $\mathrm{P}$-algebra formally smooth for the discrete topology \sref{0.19.6.1}, and consequently the $\mathrm{P}$-extension $\mathrm{A}/\mathfrak{m}^2$ of $\mathrm{K}$ by $\mathfrak{m}/\mathfrak{m}^2$ is $\mathrm{P}$-\emph{trivial}, and the characteristic homomorphism
\[
\label{0.22.2.1.1}
  \chi_{\mathrm{A}/k}:\Upsilon_{\mathrm{K}/k}\to\mathfrak{m}/\mathfrak{m}^2
  \tag{22.2.1.1}
\]
is therefore defined \sref{0.20.6.23}.
\end{env}

\begin{theorem}[22.2.2]
\label{0.22.2.2}
Under the conditions of \sref{0.22.2.1}, for $\mathrm{A}$ to be a $k$-algebra formally smooth for the $\mathfrak{m}$-preadic topology, it is necessary and sufficient that the following conditions be satisfied:
\begin{enumerate}
  \item[\rm{(i)}] \emph{The canonical homomorphism $\mathbf{S}^*_\mathrm{K}(\mathfrak{m}/\mathfrak{m}^2)\to\gr^*_\mathfrak{m}(\mathrm{A})$ is bijective.}
  \item[\rm{(ii)}] \emph{The characteristic homomorphism $\chi_{\mathrm{A}/k}:\Upsilon_{\mathrm{K}/k}\to\mathfrak{m}/\mathfrak{m}^2$ is injective.}
\end{enumerate}
This is a particular case of \sref{0.22.1.2}, where one replaces $\Lambda$ by $\mathrm{P}$, $\mathrm{A}$ by $k$, and $k$ and $\mathrm{B}$ by $\mathrm{A}$; $\mathrm{K}$ and $k$, being separable over $\mathrm{P}$, are indeed formally smooth $\mathrm{P}$-algebras \sref{0.19.6.1}.
\end{theorem}

\begin{remark}[22.2.3]
\label{0.22.2.3}
The condition (ii) of \sref{0.22.2.2} can be replaced by any one of the following:
\begin{enumerate}
  \item[a)] The canonical homomorphism $\Omega^1_k\otimes_k\mathrm{K}\to\Omega^1_\mathrm{A}\otimes_\mathrm{A}\mathrm{K}$ is injective.
  \item[b)] For every $n$, the canonical homomorphism $\Omega^1_k\otimes_k(\mathrm{A}/\mathfrak{m}^{n+1})\to\Omega^1_\mathrm{A}\otimes_\mathrm{A}(\mathrm{A}/\mathfrak{m}^{n+1})$ is invertible on the left.
  \item[c)] The canonical homomorphism $\Omega^1_k\widehat{\otimes}_k\mathrm{A}\to\widehat{\Omega}^1_\mathrm{A}$ is invertible on the left.
  \item[d)] $\mathrm{A}$ is a $k$-algebra formally smooth (for the $\mathfrak{m}$-preadic topology) \emph{relative to the prime field}.
\end{enumerate}
Indeed, it was already noted in \sref{0.22.1.1} that the injectivity of $\chi_{\mathrm{A}/k}$ is equivalent to $a$), and that the conjunction of $a$) and condition (i) of \sref{0.22.2.2} is equivalent to that of (i) and $d$).
Finally, we have also seen in \sref{0.22.1.1} that $b$) and $d$) are equivalent, and the equivalence of $b$) and $c$) follows from \sref{0.19.1.8}.
\end{remark}

\begin{env}[22.2.4]
\label{0.22.2.4}
The main application of \sref{0.22.2.2} concerns Noetherian local $k$-algebras that are \emph{complete} (the case where $\mathfrak{m}/\mathfrak{m}^2$ is a $\mathrm{K}$-vector space of \emph{finite} rank).
More precisely, let $k$ be a field, $\mathrm{K}$ an extension of $k$, $\mathrm{V}$ a $\mathrm{K}$-vector space of
\oldpage[0\textsubscript{IV-1}]{187}
finite rank; we consider the triplets $(\mathrm{A},\alpha,\beta)$, where $\mathrm{A}$ is a complete Noetherian local $k$-algebra \emph{formally smooth} with maximal ideal $\mathfrak{m}$, $\alpha:\mathrm{A}/\mathfrak{m}\xrightarrow{\sim}\mathrm{K}$ a $k$-algebra isomorphism, $\beta:\mathfrak{m}/\mathfrak{m}^2\xrightarrow{\sim}\mathrm{V}$ a di-isomorphism of vector spaces (for the isomorphism $\alpha$).
Given two such triplets $(\mathrm{A},\alpha,\beta)$, $(\mathrm{A}',\alpha',\beta')$, an \emph{equivalence} from the first to the second is a $k$-isomorphism $u:\mathrm{A}\xrightarrow{\sim}\mathrm{A}'$ such that (if $\mathfrak{m}'$ is the maximal ideal of $\mathrm{A}'$), the diagrams
\[
  \begin{tikzcd}
    \mathrm{A}/\mathfrak{m} \ar[r,"\gr^0(u)"] \ar[d,"\alpha"] & \mathrm{A}'/\mathfrak{m}' \ar[d,"\alpha'"] & \mathfrak{m}/\mathfrak{m}^2 \ar[r,"\gr^1(u)"] \ar[d,"\beta"] & \mathfrak{m}'/\mathfrak{m}'^2 \ar[d,"\beta'"]\\
    \mathrm{K} \ar[r,"1_\mathrm{K}"] & \mathrm{K} & \mathrm{V} \ar[r,"1_\mathrm{V}"] & \mathrm{V}
  \end{tikzcd}
\]
are commutative.
It is immediate that the \emph{equivalence classes} of triplets $(\mathrm{A},\alpha,\beta)$ form a \emph{set} that we shall denote $\operatorname{FL}(\mathrm{K}/k,\mathrm{V})$.
We now note that to every triplet $(\mathrm{A},\alpha,\beta)$ is associated the $\mathrm{K}$-homomorphism composition of vector spaces $\Upsilon_{\mathrm{K}/k}\xrightarrow{\chi_{\mathrm{A}/k}}\mathfrak{m}/\mathfrak{m}^2\xrightarrow{\beta}\mathrm{V}$ (where $\chi_{\mathrm{A}/k}$ is considered as a di-homomorphism relative to $\alpha^{-1}$); the preceding definition shows (by \sref{0.20.6.22.2}) that this $\mathrm{K}$-homomorphism \emph{depends only on the equivalence class} of $(\mathrm{A},\alpha,\beta)$ in $\operatorname{FL}(\mathrm{K}/k,\mathrm{V})$; in other words we have thus defined a canonical map
\[
\label{0.22.2.4.1}
  \operatorname{FL}(\mathrm{K}/k,\mathrm{V})\to\Hom_\mathrm{K}(\Upsilon_{\mathrm{K}/k},\mathrm{V}).
  \tag{22.2.4.1}
\]
\end{env}

\begin{proposition}[22.2.5]
\label{0.22.2.5}
Let $k$ be a field, $\mathrm{K}$ an extension of $k$, $\mathrm{V}$ a $\mathrm{K}$-vector space of finite rank.
The canonical map \sref{0.22.2.4.1} is a bijection from $\operatorname{FL}(\mathrm{K}/k,\mathrm{V})$ to the set of injective $\mathrm{K}$-homomorphisms from $\Upsilon_{\mathrm{K}/k}$ to $\mathrm{V}$.
\end{proposition}

\begin{proof}
(i) To show that \sref{0.22.2.4.1} is injective, it must be proved that the following is true: let $\mathrm{A}$, $\mathrm{A}'$ be two complete Noetherian local $k$-algebras formally smooth, $\mathfrak{m}$, $\mathfrak{m}'$ their respective maximal ideals, $\alpha:\mathrm{A}/\mathfrak{m}\xrightarrow{\sim}\mathrm{K}$, $\alpha':\mathrm{A}'/\mathfrak{m}'\xrightarrow{\sim}\mathrm{K}$ two $k$-isomorphisms; let us be given two $k$-isomorphisms $v^0:\mathrm{A}/\mathfrak{m}\xrightarrow{\sim}\mathrm{A}'/\mathfrak{m}'$, $v^1:\mathfrak{m}/\mathfrak{m}^2\xrightarrow{\sim}\mathfrak{m}'/\mathfrak{m}'^2$ such that the diagrams
\[
\label{0.22.2.5.1}
  \begin{tikzcd}
    \mathrm{A}/\mathfrak{m} \ar[r,"v^0"] \ar[d,"\alpha"] & \mathrm{A}'/\mathfrak{m}' \ar[dl,"\alpha'"] & \mathfrak{m}/\mathfrak{m}^2 \ar[r,"v^1"] \ar[d,"\chi_\mathrm{A}"] & \mathfrak{m}'/\mathfrak{m}'^2 \ar[dl,"\chi_{\mathrm{A}'}"]\\
    \mathrm{K} & & \Upsilon_{\mathrm{K}/k} &
  \end{tikzcd}
  \tag{22.2.5.1}
\]
are commutative.
One must prove that there exists a $k$-isomorphism $u:\mathrm{A}\xrightarrow{\sim}\mathrm{A}'$ such that $\gr^0(u)=v^0$ and $\gr^1(u)=v^1$.
\oldpage[0\textsubscript{IV-1}]{188}
We note first that since $\mathrm{K}$ is a field, the canonical homomorphism \sref{0.20.6.7.1} is \emph{bijective} \sref{0.20.6.11}, so there already exists a $k$-isomorphism $v:\mathrm{A}/\mathfrak{m}^2\xrightarrow{\sim}\mathrm{A}'/\mathfrak{m}'^2$ such that $\gr^0(v)=v^0$ and $\gr^1(v)=v^1$.
We now note that since $\mathrm{A}'$ is metrizable and complete, and $\mathrm{A}$ a $k$-algebra formally smooth, the composite homomorphism $\mathrm{A}\to\mathrm{A}/\mathfrak{m}^2\xrightarrow{v}\mathrm{A}'/\mathfrak{m}'^2$ factors as $\mathrm{A}\xrightarrow{u}\mathrm{A}'\to\mathrm{A}'/\mathfrak{m}'^2$, where $u$ is a $k$-homomorphism \sref{0.19.3.11}, and one has thus already $\gr^0(u)=v^0$ and $\gr^1(u)=v^1$.
But since $v^0$ and $v^1$ are bijective, and since $\mathrm{A}$ and $\mathrm{A}'$ are $k$-algebras formally smooth, one deduces from \sref{0.22.2.2}, (i) that $\gr(u)$ is \emph{bijective}; but since $\mathrm{A}$ and $\mathrm{A}'$ are separated and complete, this implies that $u$ itself is \emph{bijective} (Bourbaki, \emph{Alg.~comm.}, chap.~III, \textsection~2, n\textsuperscript{o}~8, cor.~3 of th.~1).

(ii) It remains to show that the image of \sref{0.22.2.4.1} is the set of \emph{injective} homomorphisms of $\Upsilon_{\mathrm{K}/k}$ in $\mathrm{V}$; one already knows by \sref{0.22.2.2}, (ii) that it is contained in this set, so it remains to prove the
\end{proof}

\begin{lemma}[22.2.5.2]
\label{0.22.2.5.2}
For every $\mathrm{K}$-homomorphism $h:\Upsilon_{\mathrm{K}/k}\to\mathrm{V}$, there exists a $k$-algebra $\mathrm{A}$ that is a local Noetherian ring, regular and complete, with maximal ideal $\mathfrak{m}$, with residue field $\mathrm{K}$, and a $\mathrm{K}$-isomorphism $\beta:\mathfrak{m}/\mathfrak{m}^2\xrightarrow{\sim}\mathrm{V}$ such that $h=\beta\circ\chi_{\mathrm{A}/k}$.
\end{lemma}

\begin{proof}
Indeed, if this lemma is proved, the fact that $\mathrm{A}$ is regular implies that condition (i) of \sref{0.22.2.2} is satisfied \sref{0.17.1.1}; if moreover $h$ is injective, \sref{0.22.2.2} will show that $\mathrm{A}$ is a $k$-algebra formally smooth whose equivalence class maps to $h$ under \sref{0.22.2.4.1} (while if $h$ is not injective, the $k$-algebra $\mathrm{A}$ of which \sref{0.22.2.5.2} proves the existence is not formally smooth).

To prove \sref{0.22.2.5.2}, let $\mathrm{B}$ be the $\mathrm{K}$-algebra $\mathbf{S}^*_\mathrm{K}(\mathrm{V})$, $\mathfrak{n}$ the augmentation ideal of the completion $\widehat{\mathrm{B}}$ of $\mathrm{B}$ for the $\mathfrak{n}$-preadic topology (such that $\mathrm{A}$ is isomorphic to a formal power series algebra $\mathrm{K}[[\mathrm{T}_1,\ldots,\mathrm{T}_n]]$ if $n=\rg_\mathrm{K}\mathrm{V}$).
By virtue of \sref{0.20.6.11} there exists on $\mathrm{A}/\mathfrak{m}^2$ (where $\mathfrak{m}=\mathfrak{n}\mathrm{A}$) a structure of $k$-extension of $\mathrm{K}$ by $\mathfrak{m}/\mathfrak{m}^2$ that is $\mathrm{P}$-\emph{trivial} (by definition, since the prime field $\mathrm{P}$ is separable), such that $h$ is the characteristic homomorphism of this extension.
If $f:k\xrightarrow{g}\mathrm{A}\to\mathrm{A}/\mathfrak{m}^2$ is the structural homomorphism, the fact that $k$ is separable over the prime field permits \sref{0.19.6.1} factoring $f$ as $k\xrightarrow{g}\mathrm{A}\to\mathrm{A}/\mathfrak{m}^2$, and the $k$-algebra $\mathrm{A}$ thus defined satisfies \sref{0.20.6.22.2}.
\end{proof}

\begin{theorem}[22.2.6]
\label{0.22.2.6}
Let $k$ be a field, $\mathrm{K}$ an extension of $k$.
For there to exist a $k$-algebra $\mathrm{A}$, which is a local Noetherian ring, with maximal ideal $\mathfrak{m}$, such that $\mathrm{A}/\mathfrak{m}=\mathrm{K}$ and such that $\mathrm{A}$ is formally smooth (for the $\mathfrak{m}$-preadic topology), it is necessary and sufficient that $\Upsilon_{\mathrm{K}/k}$ is a $\mathrm{K}$-vector space of finite rank.
The $k$-algebra $\mathrm{A}$ is then determined (up to $k$-isomorphism inducing the identity $\mathrm{K}\to\mathrm{K}$ on residue fields) by its dimension alone, which is subject to the sole condition of being $\geq\rg_\mathrm{K}\Upsilon_{\mathrm{K}/k}$.
\end{theorem}

\begin{proof}
Since $\mathfrak{m}/\mathfrak{m}^2$ is of finite rank over $\mathrm{K}$, the necessity of the condition follows from \sref{0.22.2.2}, (ii); conversely, \sref{0.22.2.5} shows that the condition is sufficient and that one can take for $\mathfrak{m}/\mathfrak{m}^2$ a $\mathrm{K}$-vector space of rank any $\geq\rg_\mathrm{K}(\Upsilon_{\mathrm{K}/k})$.
\end{proof}

\oldpage[0\textsubscript{IV-1}]{189}
In particular, the identity homomorphism of $\Upsilon_{\mathrm{K}/k}$ associates \emph{canonically} to $\mathrm{K}$ a complete Noetherian $k$-algebra, formally smooth, for which $\Upsilon_{\mathrm{K}/k}$ is isomorphic to $\mathfrak{m}/\mathfrak{m}^2$.

\begin{remark}[22.2.7]
\label{0.22.2.7}
\begin{enumerate}
  \item[\rm{(i)}] Let $\mathrm{A}$ be a local Noetherian ring, with residue field $k$; it follows from \sref{0.19.7.1} and \sref{0.19.7.2} that the determination of the $\Lambda$-algebras $\mathrm{A}$ that are complete Noetherian local rings and are formally smooth is \emph{equivalent} to the same problem where $\Lambda$ is replaced by $k$, and is thus in principle entirely solved by \sref{0.22.2.5}, which reduces the question to the structure of the \emph{residue fields} of the $k$-algebras sought, as extensions of $k$.
  \item[\rm{(ii)}] The fact that $\Upsilon_{\mathrm{K}/k}$ is of finite rank over $\mathrm{K}$ does \emph{not imply} that $\mathrm{K}$ is of ``finite radiciel multiplicity'' over $k$ (cf.~\sref{0.19.6.6}).
  For example, let $k_0$ be a perfect field of characteristic $p>0$, $k=k_0(\mathrm{T})$ the field of rational functions in one indeterminate, so that $\mathrm{K}=k^{p^{-\infty}}$ is the smallest perfect extension of $k$.
  Then $\Omega^1_k=\Omega^1_{k/k}=0$ \sref{0.21.4.4}, therefore $\Upsilon_{\mathrm{K}/k}=\Omega^1_k\otimes_k\mathrm{K}$, and since $k^p=k_0(\mathrm{T}^p)$, $\mathrm{T}$ is a $p$-base of $k$ over $k^p$, therefore \sref{0.21.4.5}, $\Omega^1_k$ is of rank $1$ over $k$, and hence $\Upsilon_{\mathrm{K}/k}$ is of rank $1$ over $\mathrm{K}$.
  But for every $n$ the residue field of the local ring $\mathrm{K}\otimes_kk^{p^{-n}}$ is isomorphic to $\mathrm{K}$, which is therefore not separable over $k^{p^{-n}}$.
\end{enumerate}
\end{remark}

\begin{proposition}[22.2.8]
\label{0.22.2.8}
Under the conditions of \sref{0.22.2.1}, suppose that the characteristic $p$ of $k$ is $>0$; let $(k_\alpha)$ be a decreasing filtering family of subfields of $k$, such that $\bigcap k_\alpha(k^p)=k^p$.
Then the condition (ii) of \sref{0.22.2.2} can be replaced by any one of the following:
\begin{enumerate}
  \item[a)] There exists $\alpha_0$ such that the canonical homomorphism
  \[
    \Omega^1_{k/k_\alpha}\otimes_k\mathrm{K}\to\Omega^1_{\mathrm{A}/k_\alpha}\otimes_\mathrm{A}\mathrm{K}
  \]
  is injective for every $\alpha\geq\alpha_0$.
  \item[b)] There exists $\alpha_0$ such that $\mathrm{A}$ is a $k$-algebra formally smooth (for the $\mathfrak{m}$-preadic topology) relative to $k_\alpha$, for every $\alpha\geq\alpha_0$.
  \item[c)] For every $\alpha$ and every integer $n\geq 0$, the canonical homomorphism
  \[
    \Omega^1_{k/k_\alpha}\otimes_k(\mathrm{A}/\mathfrak{m}^{n+1})\to\Omega^1_{\mathrm{A}/k_\alpha}\otimes_\mathrm{A}(\mathrm{A}/\mathfrak{m}^{n+1})
  \]
  is invertible on the left.
\end{enumerate}
\end{proposition}

\begin{proof}
Indeed, if $\mathrm{A}$ is a formally smooth $k$-algebra, it is so \emph{a fortiori} relative to every subfield of $k$, which implies $c$) by virtue of \sref{0.20.7.2} and \sref{0.19.1.7}; it is clear that $c$) implies $b$), and that $b$) implies $a$).
Finally, taking into account \sref{0.22.2.3}, it remains to prove that $a$) implies that the canonical homomorphism $u:\Omega^1_k\otimes_k\mathrm{K}\to\Omega^1_\mathrm{A}\otimes_\mathrm{A}\mathrm{K}$ is \emph{injective}.
\oldpage[0\textsubscript{IV-1}]{190}
But for every $\alpha$, one has a commutative diagram
\[
  \begin{tikzcd}
    \Omega^1_k\otimes_k\mathrm{K} \ar[r,"u"] & \Omega^1_\mathrm{A}\otimes_\mathrm{A}\mathrm{K}\\
    \Omega^1_{k/k_\alpha}\otimes_k\mathrm{K} \ar[u,"v_\alpha"] \ar[r,"u_\alpha"] & \Omega^1_{\mathrm{A}/k_\alpha}\otimes_\mathrm{A}\mathrm{K} \ar[u]
  \end{tikzcd}
\]
and the hypothesis that $u_\alpha$ is injective for $\alpha\geq\alpha_0$ implies that $\Ker(u)$ is contained in the intersection of the $\Ker(v_\alpha)$.
But one has seen \sref{0.21.8.3} that this intersection is zero by virtue of the hypothesis $\bigcap k_\alpha(k^p)=k^p$, which finishes the proof.
\end{proof}

\begin{proposition}[22.2.9]
\label{0.22.2.9}
Under the hypotheses of \sref{0.22.2.1}, the following conditions are equivalent:
\begin{enumerate}
  \item[a)] $\mathrm{A}/\mathfrak{m}^2$ is a $k$-extension $k$-trivial of $\mathrm{K}$ by $\mathfrak{m}/\mathfrak{m}^2$.
  \item[b)] $\chi_{\mathrm{A}/k}=0$.
  \item[c)] The canonical homomorphism \sref{0.20.5.11.2}
  \[
    \delta_{\mathrm{K}/\mathrm{A}/k}:\mathfrak{m}/\mathfrak{m}^2\to\Omega^1_{\mathrm{A}/k}\otimes_\mathrm{A}\mathrm{K}
  \]
  is injective (in other words, the sequence
\end{enumerate}
\[
\label{0.22.2.9.1}
  0\to\mathfrak{m}/\mathfrak{m}^2\to\Omega^1_{\mathrm{A}/k}\otimes_\mathrm{A}\mathrm{K}\to\Omega^1_{\mathrm{K}/k}\to 0
  \tag{22.2.9.1}
\]
\emph{is exact}).
\end{proposition}

\begin{proof}
The equivalence of $b$) and $c$) follows from the exact sequence \sref{0.20.6.22.1}; the equivalence of $c$) and $a$) follows from \sref{0.20.5.12}, this sequence being split if it is exact, since these are vector spaces.

Note that if $\mathrm{K}$ is \emph{separable} over $k$, the equivalent conditions of \sref{0.22.2.9} are satisfied, by virtue of the definition of $\chi_{\mathrm{A}/k}$ \sref{0.22.2.1} and of \sref{0.20.6.3}.
\end{proof}

\begin{corollary}[22.2.10]
\label{0.22.2.10}
Under the hypotheses of \sref{0.22.2.1}, consider a subfield $k_0$ of $k$.
The following conditions are equivalent:
\begin{enumerate}
  \item[a)] $\mathrm{A}/\mathfrak{m}^2$ is a $k_0$-extension $k_0$-trivial of $\mathrm{K}$ by $\mathfrak{m}/\mathfrak{m}^2$.
  \item[b)] The composite homomorphism $\Upsilon_{\mathrm{K}/k_0}\to\Upsilon_{\mathrm{K}/k}\xrightarrow{\chi_{\mathrm{A}/k_0}}\mathfrak{m}/\mathfrak{m}^2$ is zero.
  \item[c)] The characteristic homomorphism $\chi_{\mathrm{A}/k}:\Upsilon_{\mathrm{K}/k}\to\mathfrak{m}/\mathfrak{m}^2$ factors through $\Upsilon_{\mathrm{K}/k}\to\Upsilon_{\mathrm{K}/k/k_0}\xrightarrow{\chi'}\mathfrak{m}/\mathfrak{m}^2$.
\end{enumerate}
Moreover, when these conditions are satisfied, the homomorphism $\chi'$ is determined in a unique and canonical way, and coincides with the characteristic homomorphism of the $k_0$-extension $k_0$-trivial $\mathrm{A}/\mathfrak{m}^2$ of $\mathrm{K}$ by $\mathfrak{m}/\mathfrak{m}^2$.
\end{corollary}

\begin{proof}
Since the composition in $b$) is nothing other than $\chi_{\mathrm{A}/k_0}$, the equivalence of $a$) and $b$) follows from \sref{0.22.2.9}; on the other hand, $b$) and $c$) are equivalent by virtue of the exact sequence \sref{0.20.6.18.1}; and since the canonical homomorphism $\Upsilon_{\mathrm{K}/k}\to\Upsilon_{\mathrm{K}/k/k_0}$ is surjective, this implies the uniqueness of $\chi'$; the fact that $\chi'$ is the characteristic homomorphism of the $k_0$-extension $\mathrm{A}/\mathfrak{m}^2$ follows from \sref{0.20.6.22.2}.
\end{proof}

\begin{corollary}[22.2.11]
\label{0.22.2.11}
Under the hypotheses of \sref{0.22.2.1}, let $p$ be the characteristic exponent of $k$, and let $(k_\alpha)$ be a decreasing filtering family of subfields of $k$ such that $\bigcap k_\alpha(k^p)=k^p$.
If $\Upsilon_{\mathrm{K}/k}$ is of finite rank over $\mathrm{K}$, there exists an index $\alpha$ such that $\mathrm{A}/\mathfrak{m}^2$ is a $k_\alpha$-extension $k_\alpha$-trivial of $\mathrm{K}$ by $\mathfrak{m}/\mathfrak{m}^2$.
\end{corollary}

\begin{proof}
It is known \sref{0.21.8.3} that there exists then $\alpha$ such that the canonical homomorphism $\Upsilon_{\mathrm{K}/k}\to\Upsilon_{\mathrm{K}/k/k_\alpha}$ is bijective, so the condition $b$) of \sref{0.22.2.10} (where $k_\alpha$ replaces $k_0$) is satisfied.
\end{proof}

\subsection{Application to relations between certain local rings and their completions}
\label{subsection:0.22.3}

\oldpage[0\textsubscript{IV-1}]{191}
\begin{lemma}[22.3.1]
\label{0.22.3.1}
Let $\mathrm{A}_0$ be a semi-local Noetherian ring, $\mathrm{A}$ a finite $\mathrm{A}_0$-algebra (which is thus a semi-local Noetherian ring, cf.~Bourbaki, \emph{Alg.~comm.}, chap.~IV, \textsection~2, n\textsuperscript{o}~5, cor.~3 of prop.~9).
If $\widehat{\mathrm{A}}_0$ and $\widehat{\mathrm{A}}$ are the completions of $\mathrm{A}_0$ and $\mathrm{A}$ for their preadic topologies respectively, then, when one equips $\mathrm{A}_0$, $\mathrm{A}$ and $\widehat{\mathrm{A}}$ with discrete topologies, $\widehat{\mathrm{A}}$ is an $\mathrm{A}$-algebra formally smooth relative to $\mathrm{A}_0$ \sref{0.19.9.1}.
\end{lemma}

\begin{proof}
It is known that $\widehat{\mathrm{A}}=\mathrm{A}\otimes_{\mathrm{A}_0}\widehat{\mathrm{A}}_0$ (Bourbaki, \emph{Alg.~comm.}, chap.~IV, \textsection~2, n\textsuperscript{o}~5, cor.~3 of prop.~9 and chap.~III, \textsection~3, n\textsuperscript{o}~4, th.~1), so the proposition follows from \sref{0.19.9.3}.
\end{proof}

\begin{proposition}[22.3.2]
\label{0.22.3.2}
Let $\mathrm{A}$ be a local Noetherian regular ring, $\mathrm{K}$ its field of fractions, $p$ the characteristic of $\mathrm{K}$.
Suppose that one of the following hypotheses is satisfied:
\begin{enumerate}
  \item[\rm{(i)}] $p=0$.
  \item[\rm{(ii)}] $p>0$ and there exists a filtering decreasing family $(\mathrm{A}_\alpha)$ of Noetherian subrings of $\mathrm{A}$, such that $\mathrm{A}$ is a finite $\mathrm{A}_\alpha$-algebra for every $\alpha$, such that $\mathrm{A}^{p^{h(\alpha)}}\subset\mathrm{A}_\alpha$ for a suitable integer $h(\alpha)$, and such that, if $\mathrm{K}_\alpha$ denotes the field of fractions of $\mathrm{A}_\alpha$, we have $\bigcap\mathrm{K}_\alpha(\mathrm{K}^p)=\mathrm{K}^p$.
\end{enumerate}

Let then $\mathrm{B}$ be a finite integral $\mathrm{A}$-algebra containing $\widehat{\mathrm{A}}$, $\mathfrak{n}$ a prime ideal of $\mathrm{B}$, $\mathrm{C}$ the local ring $\mathrm{B}_\mathfrak{n}$, $\mathfrak{q}$ a prime ideal of the completion $\widehat{\mathrm{C}}$, such that $\mathfrak{q}\cap\mathrm{C}=0$, so that the local ring $\widehat{\mathrm{C}}_\mathfrak{q}$ is an algebra over the field of fractions $\mathrm{L}$ of $\mathrm{B}$.
Then $\widehat{\mathrm{C}}_\mathfrak{q}$ is an $\mathrm{L}$-algebra formally smooth for its $\mathfrak{q}$-adic topology, and is therefore a geometrically regular ring over $\mathrm{L}$ \sref{0.19.6.5}.
\end{proposition}

\begin{proof}
We distinguish two cases.

\emph{I)} Suppose that $\mathfrak{n}$ contains the maximal ideal $\mathfrak{m}$ of $\mathrm{A}$, so that $\mathfrak{n}\cap\mathrm{A}=\mathfrak{m}$.
Then $\mathfrak{n}$ is maximal.
If $\mathfrak{r}$ is the radical of the semi-local ring $\mathrm{B}$, $\widehat{\mathrm{C}}$ is the separated completion of $\mathrm{B}$ for the $\mathfrak{n}$-preadic topology, and a component of the semi-local ring $\widehat{\mathrm{B}}$, the completion of $\mathrm{B}$ for the $\mathfrak{r}$-preadic topology (Bourbaki, \emph{Alg.~comm.}, chap.~III, \textsection~3, n\textsuperscript{o}~4, prop.~8);
\oldpage[0\textsubscript{IV-1}]{192}
one can thus write $\widehat{\mathrm{C}}_\mathfrak{q}=\widehat{\mathrm{B}}_{\mathfrak{q}'}$, where $\mathfrak{q}'$ is the prime ideal of $\widehat{\mathrm{B}}$ that is the inverse image of $\mathfrak{q}$.
Let us now note that $\mathrm{B}$ is a subring of $\widehat{\mathrm{B}}$, and the hypothesis $\mathfrak{q}\cap\mathrm{C}=0$ implies $\mathfrak{q}'\cap\mathrm{B}=0$, so $\widehat{\mathrm{B}}_{\mathfrak{q}'}$ is also a localization of $\widehat{\mathrm{B}}\otimes_\mathrm{B}\mathrm{L}$ at one of its prime ideals $\mathfrak{q}''$, of which $\mathfrak{q}'$ is the inverse image in $\widehat{\mathrm{B}}$;
moreover we have $\widehat{\mathrm{B}}=\widehat{\mathrm{A}}\otimes_\mathrm{A}\mathrm{B}$ (Bourbaki, \emph{Alg.~comm.}, chap.~IV, \textsection~2, n\textsuperscript{o}~5, cor. of prop.~9 and chap.~III, \textsection~3, n\textsuperscript{o}~4, th.~1), so $\widehat{\mathrm{B}}\otimes_\mathrm{B}\mathrm{L}=\widehat{\mathrm{A}}\otimes_\mathrm{A}\mathrm{L}$; finally, if $\mathfrak{p}$ is the prime ideal of $\widehat{\mathrm{A}}$ that is the inverse image of $\mathfrak{q}''$, the fact that $\mathrm{A}\to\mathrm{B}$ is injective and the commutativity of the diagram
\[
  \begin{tikzcd}
    \mathrm{B} \ar[r] & \widehat{\mathrm{B}}\\
    \mathrm{A} \ar[u] \ar[r] & \widehat{\mathrm{A}} \ar[u]
  \end{tikzcd}
\]
imply that $\mathfrak{p}\cap\mathrm{A}=0$, so $(\widehat{\mathrm{A}}\otimes_\mathrm{A}\mathrm{L})_{\mathfrak{q}''}$ is also a localization of the ring $\widehat{\mathrm{A}}_\mathfrak{p}\otimes_\mathrm{A}\mathrm{L}=(\widehat{\mathrm{A}}_\mathfrak{p}\otimes_\mathrm{A}\mathrm{K})\otimes_\mathrm{K}\mathrm{L}$.
To show that this local ring is an $\mathrm{L}$-algebra formally smooth, it suffices therefore to prove that $\widehat{\mathrm{A}}_\mathfrak{p}\otimes_\mathrm{A}\mathrm{K}$ is a $\mathrm{K}$-algebra \emph{formally smooth}.
First of all, as $\mathrm{A}$ is regular, so is $\widehat{\mathrm{A}}$ \sref{0.17.3.2}, and condition (i) of \sref{0.22.2.2} is satisfied \sref{0.17.1.1}.
It therefore suffices (when $p>0$) to verify condition $b$) of \sref{0.22.2.8}.
But \sref{0.22.3.1} shows that, for every $\alpha$, $\widehat{\mathrm{A}}$ is an $\mathrm{A}$-algebra formally smooth relative to $\mathrm{A}_\alpha$ (for the discrete topologies); hence $\widehat{\mathrm{A}}\otimes_{\mathrm{A}_\alpha}\mathrm{K}_\alpha=\widehat{\mathrm{A}}\otimes_\mathrm{A}(\mathrm{A}\otimes_{\mathrm{A}_\alpha}\mathrm{K}_\alpha)$ is an $(\mathrm{A}\otimes_{\mathrm{A}_\alpha}\mathrm{K}_\alpha)$-algebra formally smooth relative to $\mathrm{K}_\alpha$, for the discrete topologies \sref{0.19.9.2}, (iii).
Since $\mathrm{A}$ is an $\mathrm{A}_\alpha$-algebra \emph{finite}, we have $\mathrm{A}\otimes_{\mathrm{A}_\alpha}\mathrm{K}_\alpha=\mathrm{K}$; as $\widehat{\mathrm{A}}_\mathfrak{p}$ is a ring of fractions of $\widehat{\mathrm{A}}\otimes_\mathrm{A}\mathrm{K}$, we deduce \sref{0.19.9.2}, (iv) that $\widehat{\mathrm{A}}_\mathfrak{p}$ is a $\mathrm{K}$-algebra formally smooth \emph{relative} to $\mathrm{K}_\alpha$, for the discrete topology, and \emph{a fortiori} for its $\mathfrak{p}$-adic topology \sref{0.19.3.8} and \sref{0.19.9.5}.
But this is precisely condition $b$) of \sref{0.22.2.8} in this case.
When $p=0$, the residue field of $\widehat{\mathrm{A}}_\mathfrak{p}$ is separable over $\mathrm{K}$, and the proposition is proved in this case by virtue of \sref{0.19.6.4}.

\emph{II)} General case.
Let $\mathfrak{s}$ be the prime ideal $\mathfrak{n}\cap\mathrm{A}$ of $\mathrm{A}$, and set $\mathrm{S}=\mathrm{A}-\mathfrak{s}$, so that $\mathrm{S}^{-1}\mathrm{A}=\mathrm{A}_\mathfrak{s}$, and $\mathrm{C}=(\mathrm{S}^{-1}\mathrm{B})_{\mathfrak{s}\cdot\mathfrak{n}}$, where this time $\mathrm{S}^{-1}\mathfrak{n}$ contains the maximal ideal $\mathfrak{s}\mathrm{A}_\mathfrak{s}$ of $\mathrm{A}_\mathfrak{s}$.
As $\mathrm{A}_\mathfrak{s}$ is regular \sref{0.17.3.2} and $\mathrm{S}^{-1}\mathrm{B}$ a finite $\mathrm{A}_\mathfrak{s}$-algebra integral containing $\mathrm{A}_\mathfrak{s}$, it suffices therefore to verify condition (ii) for $\mathrm{A}_\mathfrak{s}$ when $p>0$: but setting $\mathfrak{s}_\alpha=\mathfrak{s}\cap\mathrm{A}_\alpha$ and considering the local ring $(\mathrm{A}_\alpha)_{\mathfrak{s}_\alpha}$.
For every $t\notin\mathfrak{s}$, one has by hypothesis, setting $q=p^{h(\alpha)}$, $t^q\in\mathrm{A}_\alpha$ and $t^q\notin\mathfrak{s}_\alpha$; if one sets $\mathrm{S}_\alpha=\mathrm{A}_\alpha-\mathfrak{s}_\alpha$, one sees therefore that $\mathrm{A}_\mathfrak{s}$ is a finite $(\mathrm{A}_\alpha)_{\mathfrak{s}_\alpha}$-algebra; moreover, $\mathrm{K}_\alpha$ is also the field of fractions of $(\mathrm{A}_\alpha)_{\mathfrak{s}_\alpha}$, which finishes the proof.
\end{proof}

\begin{theorem}[22.3.3]
\label{0.22.3.3}
Let $\mathrm{A}$ be a local Noetherian complete ring, $\mathfrak{p}$ a prime ideal of $\mathrm{A}$, $\mathrm{B}$ the localized ring $\mathrm{A}_\mathfrak{p}$, $\mathfrak{q}$ a prime ideal of the completion $\widehat{\mathrm{B}}$, $\mathfrak{r}=\mathfrak{q}\cap\mathrm{B}$.
Then the localized ring $\widehat{\mathrm{B}}_\mathfrak{q}$ is a $\mathrm{B}_\mathfrak{r}$-algebra formally smooth for the preadic topologies.
\end{theorem}

\begin{proof}
The prime ideal $\mathfrak{r}$ of $\mathrm{B}$ is of the form $\mathfrak{n}_\mathfrak{p}$, where $\mathfrak{n}$ is a prime ideal of $\mathrm{A}$, and $\mathrm{B}_\mathfrak{r}=\mathrm{A}_\mathfrak{n}$; let $k$ be the residue field of $\mathrm{B}_\mathfrak{r}$, which is also the field of fractions of $\mathrm{A}/\mathfrak{n}$.
\oldpage[0\textsubscript{IV-1}]{193}
As $\widehat{\mathrm{B}}$ is a flat $\mathrm{B}_\mathfrak{r}$-module $(\mathbf{0}_\mathrm{I},~7.3.5)$, $\widehat{\mathrm{B}}_\mathfrak{q}$ is also a flat $\mathrm{B}_\mathfrak{r}$-module, and to prove that $\widehat{\mathrm{B}}_\mathfrak{q}$ is a $\mathrm{B}_\mathfrak{r}$-algebra formally smooth, it suffices to prove that $\widehat{\mathrm{B}}_\mathfrak{q}\otimes_{\mathrm{B}_\mathfrak{r}}k=\widehat{\mathrm{B}}_\mathfrak{q}/\mathfrak{r}\widehat{\mathrm{B}}_\mathfrak{q}$ is a $k$-algebra formally smooth.
But $\widehat{\mathrm{B}}_\mathfrak{q}/\mathfrak{r}\widehat{\mathrm{B}}_\mathfrak{q}=(\widehat{\mathrm{B}}/\mathfrak{r}\widehat{\mathrm{B}})_{\mathfrak{q}'}$, where $\mathfrak{q}'$ is the canonical image of $\mathfrak{q}$ in $\widehat{\mathrm{B}}/\mathfrak{r}\widehat{\mathrm{B}}$; one has by definition $\mathfrak{q}'\cap(\mathrm{B}/\mathfrak{r})=0$ and $\mathrm{B}/\mathfrak{r}=\mathrm{A}_{\mathfrak{p}'}$, where $\mathfrak{p}'=\mathfrak{p}/\mathfrak{n}$; moreover $\widehat{\mathrm{B}}/\mathfrak{r}\widehat{\mathrm{B}}$ is the completion of $\mathrm{B}/\mathfrak{r}$.

It follows from \sref{0.19.8.9} that there exists a subring $\mathrm{A}_0$ of $\mathrm{A}$ that is a power series ring over a field of characteristic $p>0$ or over a Cohen ring of characteristic $0$; it is known that $\mathrm{A}$ is such that $\mathrm{A}$ is a finite $\mathrm{A}_0$-algebra.
Now, we know that the ring $\mathrm{A}_0$ is \emph{regular} \sref{0.17.3.8}, and it satisfies one of the hypotheses (i), (ii) of \sref{0.22.3.2}, by virtue of \sref{0.21.8.8}; it suffices therefore to apply \sref{0.22.3.2}, replacing $\mathrm{B}$ by $\mathrm{A}$ and $\mathrm{A}$ by $\mathrm{A}_0$.
\end{proof}

\begin{corollary}[22.3.4]
\label{0.22.3.4}
Let $\mathrm{A}$ be a local Noetherian complete ring, integral, $\mathrm{K}$ its field of fractions, $\mathfrak{p}$ a prime ideal of $\mathrm{A}$, $\mathrm{B}$ the localized ring $\mathrm{A}_\mathfrak{p}$.
Then, for every prime ideal $\mathfrak{q}$ of the completion $\widehat{\mathrm{B}}$, such that $\mathfrak{q}\cap\mathrm{B}=0$, the local ring $\widehat{\mathrm{B}}_\mathfrak{q}$ is a $\mathrm{K}$-algebra formally smooth for its $\mathfrak{q}$-adic topology, and is therefore a geometrically regular ring over $\mathrm{K}$.
\end{corollary}

\begin{proof}
Indeed, it follows from \sref{0.19.8.9} that there exists a subring $\mathrm{A}_0$ of $\mathrm{A}$ that is a power series ring over a field of characteristic $p>0$ or over a Cohen ring of characteristic $0$; it is known that the ring $\mathrm{A}_0$ is \emph{regular} \sref{0.17.3.8}, and it satisfies one of hypotheses (i), (ii) of \sref{0.22.3.2}, by virtue of \sref{0.21.8.8}; it suffices therefore to apply \sref{0.22.3.2}, replacing $\mathrm{B}$ by $\mathrm{A}$ and $\mathrm{A}$ by $\mathrm{A}_0$.
\end{proof}

\subsection{Preliminary results on finite extensions of local rings whose maximal ideal has square zero}
\label{subsection:0.22.4}
\label{0.22.4}

\begin{proposition}[22.4.1]
\label{0.22.4.1}
Let $\mathrm{A}$ be a local ring whose maximal ideal $\mathfrak{m}$ is of square zero, $\mathrm{K}=\mathrm{A}/\mathfrak{m}$ its residue field; let us denote by $\mathrm{V}$ the ideal $\mathfrak{m}$ considered as a $\mathrm{K}$-vector space.
Let $\mathrm{T}_i$ $(1\leq i\leq r)$ be indeterminates, and for each $i$, let $\mathrm{F}_i$ be a unitary polynomial of $\mathrm{A}[\mathrm{T}_i]$, of degree $d_i$; let $\mathrm{A}'$ be the quotient of the polynomial ring $\mathrm{A}[\mathrm{T}_1,\ldots,\mathrm{T}_r]$ by the ideal $\mathfrak{b}$ generated by the $r$ polynomials $\mathrm{F}_i$.
Suppose that $\mathrm{A}'$ is a local ring, and let $\mathfrak{m}'$ and $\mathrm{K}'=\mathrm{A}'/\mathfrak{m}'$ be its maximal ideal and residue field.
Then:
\begin{enumerate}
  \item[\rm{(i)}] If one sets $d=\prod_{i=1}^r d_i$, $\mathrm{A}'$ is a free $\mathrm{A}$-module of rank $d$, and we have $[\mathrm{K}':\mathrm{K}]\leq d$.
  \item[\rm{(ii)}] For $[\mathrm{K}':\mathrm{K}]=d$, it is necessary and sufficient that $\mathrm{A}'\otimes_\mathrm{A}\mathrm{K}=\mathrm{A}'/\mathfrak{m}\mathrm{A}'$ be a field (necessarily isomorphic to $\mathrm{K}'$), in other words, that $\mathfrak{m}'=\mathfrak{m}\mathrm{A}'$.
  In this case, $\mathfrak{m}'^2=0$, and if $\mathrm{V}'$ denotes the ideal $\mathfrak{m}'$ considered as a $\mathrm{K}'$-vector space, the canonical homomorphism $\mathrm{V}\otimes_\mathrm{K}\mathrm{K}'\to\mathrm{V}'$ is bijective.
\end{enumerate}
\end{proposition}

\begin{proof}
Without assuming $\mathrm{A}'$ local, it is evident by Euclidean division and recurrence on $r$, that the monomials $\mathrm{M}_\nu=t_1^{\nu_1}t_2^{\nu_2}\cdots t_r^{\nu_r}$, where $0\leq\nu_i\leq d_i-1$, form a base of $\mathrm{A}'$ as an $\mathrm{A}$-module.
If $\mathrm{A}'$ is assumed local, $\mathrm{K}'$ is a quotient of $\mathrm{A}'\otimes_\mathrm{A}\mathrm{K}$, so $[\mathrm{K}':\mathrm{K}]\leq[\mathrm{A}'\otimes_\mathrm{A}\mathrm{K}:\mathrm{K}]=d$, and there cannot be equality unless $\mathfrak{m}'=\mathfrak{m}\mathrm{A}'$, which proves the first assertion of (ii).
On the other hand, it is clear that when $\mathfrak{m}'=\mathfrak{m}\mathrm{A}'$, we have $\mathfrak{m}'^2=0$, and $\mathrm{V}\otimes_\mathrm{K}\mathrm{K}'=\mathfrak{m}\otimes_\mathrm{A}\mathrm{A}'=\mathfrak{m}\mathrm{A}'=\mathrm{V}'$ since $\mathrm{A}'$ is a free $\mathrm{A}$-module.
\end{proof}

\oldpage[0\textsubscript{IV-1}]{194}

\begin{remark}[22.4.2]
\label{0.22.4.2}
When $\mathrm{A}$ is Artinian (that is to say $\rg_\mathrm{K}(\mathrm{V})$ is finite), the hypothesis that $\mathrm{A}'$ is local always implies $\rg_{\mathrm{K}'}(\mathfrak{m}'/\mathfrak{m}'^2)\geq\rg_\mathrm{K}(\mathrm{V})$, as we shall see further on \sref{0.22.5.2}.
Note that these two ranks can be equal without having $\mathfrak{m}'=\mathfrak{m}\mathrm{A}'$.
\end{remark}

\begin{proposition}[22.4.3]
\label{0.22.4.3}
With the hypotheses and notations of \sref{0.22.4.1}, suppose that the $\mathrm{F}_i$ are of the form
\[
\label{0.22.4.3.1}
  \mathrm{F}_i(\mathrm{T}_i)=\mathrm{T}_i^{d_i}+\sum_{1\leq k\leq d_i}c_{ik}\mathrm{T}_i^{d_i-k}
  \tag{22.4.3.1}
\]
where $d_i\geq 2$ for every $i$ and $c_{ik}\in\mathfrak{m}$ (``Eisenstein polynomials'').
Then:
\begin{enumerate}
  \item[\rm{(i)}] $\mathrm{A}'$ is a local ring, and its residue field $\mathrm{K}'=\mathrm{A}'/\mathfrak{m}'$ is isomorphic to $\mathrm{K}$.
  \item[\rm{(ii)}] The kernel of the canonical homomorphism $\mathrm{V}\to\mathrm{V}'=\mathfrak{m}'/\mathfrak{m}'^2$ is the ideal $\mathfrak{n}$ of $\mathrm{A}$ generated by the elements $\xi_i=c_{i,d_i}$, and its cokernel has as a basis over $\mathrm{K}'$ the images of the $\mathrm{T}_i$.
  In particular,
  \begin{equation}
  \label{0.22.4.3.2}
    \rg_{\mathrm{K}'}(\mathrm{V}')=\rg_\mathrm{K}(\mathrm{V})+r-r',
    \tag{22.4.3.2}
  \end{equation}
  where $r'=\rg_\mathrm{K}(\mathfrak{n})$.
  When $\rg_\mathrm{K}(\mathrm{V})$ is finite, equivalently when $\mathrm{A}$ is Artinian, equality $\rg_{\mathrm{K}'}(\mathrm{V}')=\rg_\mathrm{K}(\mathrm{V})$ holds if and only if the $\xi_i$ are linearly independent over $\mathrm{K}$.
\end{enumerate}
\end{proposition}

\begin{proof}
(i) It is clear that $\mathrm{A}'\otimes_\mathrm{A}\mathrm{K}=\mathrm{A}'/\mathfrak{m}\mathrm{A}'$ is isomorphic to $\mathrm{K}[\mathrm{T}_1,\ldots,\mathrm{T}_r]/\mathfrak{b}'$, where $\mathfrak{b}'$ is the ideal generated by the $r$ polynomials $\mathrm{T}_i^{d_i}$; this is therefore a local ring of residue field $\mathrm{K}$, and it is the same for $\mathrm{A}'$, from which (i) follows.

(ii) Let $t_i$ $(1\leq i\leq r)$ be the class of $\mathrm{T}_i$ mod.~$\mathfrak{b}$; it follows from (i) that the maximal ideal $\mathfrak{m}'$ of $\mathrm{A}'$ is given by
\[
  \mathfrak{m}'=\mathfrak{m}\mathrm{A}'+\sum_i t_i\mathrm{A}'
\]
from which, taking into account that $\mathfrak{m}^2=0$,
\[
  \mathfrak{m}'^2=\sum_{i,j}t_it_j\mathrm{A}'+\sum_i t_i\mathfrak{m}\mathrm{A}'.
\]
One deduces that $\mathrm{A}'/\mathfrak{m}'^2$ is isomorphic to the ring $\mathrm{A}[\mathrm{T}_1,\ldots,\mathrm{T}_r]/\mathfrak{c}$, where $\mathfrak{c}$ is the ideal generated by the elements
\[
  \mathrm{F}_i\ (1\leq i\leq r),\quad \mathrm{T}_i\mathrm{T}_j\ (1\leq i\leq r,\ 1\leq j\leq r),\quad x\mathrm{T}_i\ (1\leq i\leq r,\ x\in\mathfrak{m}).
\]

Since $d_i\geq 2$ and $\mathrm{T}_i^{d_i}\in\mathfrak{c}$, the hypotheses show that $\mathfrak{c}$ is also generated by
\[
  \xi_i\ (1\leq i\leq r),\quad \mathrm{T}_i\mathrm{T}_j\ (1\leq i\leq r,\ 1\leq j\leq r),\quad \mathrm{T}_i x\ (1\leq i\leq r,\ x\in\mathfrak{m}).
\]

Let $\mathfrak{c}_1=\sum\mathrm{T}_i\mathrm{T}_j\mathrm{A}[\mathrm{T}_1,\ldots,\mathrm{T}_r]$, $\mathfrak{c}_2=\mathfrak{c}_1+\sum\mathfrak{m}\mathrm{T}_i\mathrm{A}[\mathrm{T}_1,\ldots,\mathrm{T}_r]$, then $\mathfrak{c}=\mathfrak{c}_2+\sum\xi_i\mathrm{A}[\mathrm{T}_1,\ldots,\mathrm{T}_r]$.
It is clear that $\mathrm{A}[\mathrm{T}_1,\ldots,\mathrm{T}_r]/\mathfrak{c}_1$ is the $\mathrm{A}$-module direct sum of $\mathrm{A}$ and of $\mathrm{A}t_i^{(1)}$, where $t_i^{(1)}$ is the class of $\mathrm{T}_i$ mod.~$\mathfrak{c}_1$.
One deduces that $\mathrm{A}[\mathrm{T}_1,\ldots,\mathrm{T}_r]/\mathfrak{c}_2$ is direct sum of $\mathrm{A}$ and of $\mathrm{K}t_i^{(2)}$, where $t_i^{(2)}$ is the class of $\mathrm{T}_i$ mod.~$\mathfrak{c}_2$.
Finally $\mathrm{A}[\mathrm{T}_1,\ldots,\mathrm{T}_r]/\mathfrak{c}$ is direct sum of $\mathrm{A}/\mathfrak{n}$ and of $\mathrm{K}t_i'$, where $t_i'$ is the class of $\mathrm{T}_i$ mod.~$\mathfrak{c}$, and $\mathfrak{n}$ is the ideal of $\mathrm{A}$ generated by the $\xi_i$; in the identification of $\mathrm{A}'/\mathfrak{m}'^2$ to $\mathrm{A}[\mathrm{T}_1,\ldots,\mathrm{T}_r]/\mathfrak{c}$, the image of $\mathrm{V}$ identifies with $\mathfrak{m}/\mathfrak{n}$, and this proves (ii).
\end{proof}

\begin{proposition}[22.4.4]
\label{0.22.4.4}
With the hypotheses and notations of \sref{0.22.4.1}, suppose that $\mathrm{K}$ is of characteristic $p>0$ (hence $p.1\in\mathfrak{m}$ in $\mathrm{A}$) and that the $\mathrm{F}_i$ are of the form
\[
\label{0.22.4.4.1}
  \mathrm{F}_i(\mathrm{T}_i)=\mathrm{T}_i^p+p\sum_{k=1}^{p-1}c_{ik}\mathrm{T}_i^{p-k}-f_i
  \tag{22.4.4.1}
\]
with $c_{ik}\in\mathrm{A}$ and $f_i\in\mathrm{A}$ for every $i$ and every $k$ such that $1\leq k\leq p-1$; suppose furthermore that if $a_i$ is the class of $f_i$ mod.~$\mathfrak{m}$.
Suppose that the family $(a_i)_{1\leq i\leq r}$ is $p$-free over $\mathrm{K}^p$ (``absolutely $p$-free'').
Then:
\begin{enumerate}
  \item[\rm{(i)}] $\mathrm{A}'$ is a local ring, with maximal ideal $\mathfrak{m}\mathrm{A}'$, whose residue field $\mathrm{K}'$ is isomorphic to the extension of $\mathrm{K}$ obtained by adjunction of the $a_i^{1/p}$ $(1\leq i\leq r)$.
  \item[\rm{(ii)}] \emph{The elements $d_{\mathrm{A}'}f_i\otimes 1_{\mathrm{K}'}$ form a base of the kernel} $\mathrm{N}=\Upsilon^{\mathrm{K}'}_{\mathrm{A}'/\mathrm{A}}$ \emph{of the canonical homomorphism} $\Omega^1_\mathrm{A}\otimes_\mathrm{A}\mathrm{K}'\to\Omega^1_{\mathrm{A}'}\otimes_{\mathrm{A}'}\mathrm{K}'$.
\end{enumerate}
\end{proposition}

\begin{proof}
\oldpage[0\textsubscript{IV-1}]{195}
(i) This time $\mathrm{A}'\otimes_\mathrm{A}\mathrm{K}=\mathrm{A}'/\mathfrak{m}\mathrm{A}'$ is isomorphic to $\mathrm{K}[\mathrm{T}_1,\ldots,\mathrm{T}_r]/\mathfrak{b}'$, where $\mathfrak{b}'$ is the ideal generated by the polynomials $\mathrm{F}_i=\mathrm{T}_i^p-b_i$ $(1\leq i\leq r)$; one deduces that $\mathrm{A}_1'$ is isomorphic to the quotient of the polynomial ring $\mathrm{A}_1[\mathrm{T}_1,\ldots,\mathrm{T}_r]$ by the ideal generated by the $\mathrm{F}_i$.
The hypothesis made on the $a_i$ implies that this quotient ring is a field \sref{0.21.1.8}, from which the assertions of (i) follow.

Before proving (ii), we shall establish the

\begin{lemma}[22.4.4.2]
\label{0.22.4.4.2}
Let $\Lambda$ be a ring, $\mathrm{A}$ a $\Lambda$-algebra, $\mathfrak{J}$ an ideal of $\mathrm{A}$ such that the ring $\mathrm{C}=\mathrm{A}/\mathfrak{J}$ is of characteristic $p>0$ (which is equivalent to saying that $p.1\in\mathfrak{J}$ in $\mathrm{A}$).
Let $\pi$ be the canonical image of $p.1$ in $\mathfrak{J}/\mathfrak{J}^2$.
If $\mathrm{C}$ is a $(\Lambda/p\Lambda)$-algebra formally smooth for the discrete topologies, one has a canonical exact sequence
\[
\label{0.22.4.4.3}
  0\to(\mathfrak{J}/\mathfrak{J}^2)/\mathrm{C}.\pi\to\Omega^1_{\Lambda/\Lambda}\otimes_\mathrm{A}\mathrm{C}\to\Omega^1_{\mathrm{C}/\Lambda}\to 0.
  \tag{22.4.4.3}
\]
\end{lemma}

\begin{proof}[Proof of the lemma]
Indeed, set $\mathrm{A}'=\mathrm{A}/p\mathrm{A}$, $\mathfrak{J}'=\mathfrak{J}/p\mathrm{A}$, so that $\mathrm{C}=\mathrm{A}'/\mathfrak{J}'$.
As $\mathrm{A}'=\mathrm{A}\otimes_\mathrm{A}(\Lambda/p\Lambda)$, we have $\Omega^1_{\Lambda/\Lambda}\otimes_\mathrm{A}\mathrm{A}'=\Omega^1_{\mathrm{A}'/\Lambda'}$ where $\Lambda'$ denotes $\Lambda/p\Lambda$ \sref{0.20.5.5}.
Applying the exact sequence \sref{0.20.5.14.1} to the $\Lambda'$-algebra formally smooth $\mathrm{C}=\mathrm{A}'/\mathfrak{J}'$, one obtains the exact sequence
\[
  0\to\mathfrak{J}'/\mathfrak{J}'^2\to\Omega^1_{\mathrm{A}'/\Lambda'}\otimes_{\mathrm{A}'}\mathrm{C}\to\Omega^1_{\mathrm{C}/\Lambda}\to 0
\]
\end{proof}

\begin{env}[22.4.4.4]
\label{0.22.4.4.4}
(ii) As $\mathfrak{m}^2=0$, we shall also denote by $\mathrm{V}'$ the ideal $\mathfrak{m}'$ considered as a $\mathrm{K}'$-vector space, and we shall set
\[
  \mathrm{V}_0=\mathrm{V}/\mathrm{K}.p=\mathfrak{m}/(\mathfrak{m}^2+p\mathrm{A}),\qquad \mathrm{V}_0'=\mathrm{V}'/\mathrm{K}'.p=\mathfrak{m}'/(\mathfrak{m}'^2+p\mathrm{A}').
\]

Applying the exact sequence \sref{0.22.4.4.3} in the case where $\Lambda=\mathbf{Z}$, to the $\mathbf{Z}$-algebra $\mathrm{A}$ (resp.~$\mathrm{A}'$) and to its ideal $\mathfrak{m}$ (resp.~$\mathfrak{m}'$), it follows, since $\mathrm{K}$ (resp.~$\mathrm{K}'$) is separable over the prime field $\mathbf{Z}/p\mathbf{Z}$, hence a $(\mathbf{Z}/p\mathbf{Z})$-algebra formally smooth \sref{0.19.6.1}, the exact sequences
\[
\label{0.22.4.4.5}
  0\to\mathrm{V}_0\to\Omega^1_\mathrm{A}\otimes_\mathrm{A}\mathrm{K}\to\Omega^1_\mathrm{K}\to 0,\qquad
  0\to\mathrm{V}_0'\to\Omega^1_{\mathrm{A}'}\otimes_{\mathrm{A}'}\mathrm{K}'\to\Omega^1_{\mathrm{K}'}\to 0.
  \tag{22.4.4.5}
\]

\oldpage[0\textsubscript{IV-1}]{196}
Consider then the diagram \sref{0.22.4.4.6} in which the two middle columns are the second exact sequence \sref{0.22.4.4.5} and the first exact sequence \sref{0.22.4.4.5} tensorized by $\mathrm{K}'$; as the first exact sequence \sref{0.22.4.4.5} is formed of $\mathrm{K}$-vector spaces, the two middle columns of \sref{0.22.4.4.6} are exact; moreover, the diagram formed by these columns and the horizontal arrows connecting them is commutative, by virtue of the proof of \sref{0.22.4.4.2} and the commutativity of the diagram of \sref{0.20.5.11.3}.
The last line of \sref{0.22.4.4.6} is the exact sequence obtained from \sref{0.20.6.1.1} by replacing $\mathrm{A}$, $\mathrm{B}$, $\mathrm{C}$ by $\mathbf{Z}$, $\mathrm{K}$ and $\mathrm{K}'$; the middle line is obtained by tensorization with $\mathrm{K}'$ from the $\mathrm{K}$-module exact sequence deduced from \sref{0.20.6.1.1} by replacing $\mathrm{A}$, $\mathrm{B}$, $\mathrm{C}$ by $\mathbf{Z}$, $\mathrm{A}$, $\mathrm{K}$ respectively.
The diagram formed by these two lines and the vertical arrows connecting them is commutative by virtue of the commutativity of \sref{0.20.6.4.3}.
We note that we are in the conditions of application of \sref{0.22.1.1}, (ii), so the upper line of the diagram \sref{0.22.4.4.6} is \emph{exact}; the extreme columns of the diagram

\[
\label{0.22.4.4.6}
  \xymatrix{
    & 0 \ar[d] & 0 \ar[d] & & \\
    0 \ar[r] & \mathrm{V}_0\otimes_\mathrm{K}\mathrm{K}' \ar[r] \ar[d] & \mathrm{V}_0' \ar[r] \ar[d] & 0 & \\
    0 \ar[r] & \mathrm{N} \ar[r] \ar[d] & \Omega^1_\mathrm{A}\otimes_\mathrm{A}\mathrm{K}' \ar[r] \ar[d] & \Omega^1_{\mathrm{A}'}\otimes_{\mathrm{A}'}\mathrm{K}' \ar[r] \ar[d] & 0\\
    0 \ar[r] & \Upsilon_{\mathrm{K}'/\mathrm{K}} \ar[r] \ar[d] & \Omega^1_\mathrm{K}\otimes_\mathrm{K}\mathrm{K}' \ar[r] \ar[d] & \Omega^1_{\mathrm{K}'} \ar[r] \ar[d] & 0\\
    & 0 & 0 & 0 &
  }
  \tag{22.4.4.6}
\]

are therefore formed respectively of the \emph{kernels} and the \emph{cokernels} of the horizontal arrows of the middle lines, which finishes proving that the diagram is \emph{commutative}.
Furthermore:
\end{env}

\begin{lemma}[22.4.4.7]
\label{0.22.4.4.7}
The homomorphisms
\[
  \mathrm{N}\to\Upsilon_{\mathrm{K}'/\mathrm{K}} \quad\text{and}\quad \Omega^1_{\mathrm{A}'}\otimes_{\mathrm{A}'}\mathrm{K}'\to\Omega^1_{\mathrm{K}'/\mathrm{K}}
\]
of diagram \sref{0.22.4.4.6} are bijective.
\end{lemma}

\begin{proof}
Indeed, the snake diagram (Bourbaki, \emph{Alg.~comm.}, chap.~I, \textsection~1, n\textsuperscript{o}~4, prop.~2) applied to the two middle columns of \sref{0.22.4.4.6} gives an exact sequence
\[
  0\to\mathrm{N}\to\Upsilon_{\mathrm{K}'/\mathrm{K}}\to 0\to\Omega^1_{\mathrm{A}'}\otimes_{\mathrm{A}'}\mathrm{K}'\to\Omega^1_{\mathrm{K}'/\mathrm{K}}\to 0.
\]
But, if $t_i$ is the canonical image of $\mathrm{T}_i$ in $\mathrm{A}'$, the relation $\mathrm{F}_i(t_i)=0$, together with the fact that $p.1\in\mathfrak{m}\mathrm{A}'\subset\mathfrak{m}'$, shows immediately that $d_{\mathrm{A}'}f_i\in\mathfrak{m}'\Omega^1_{\mathrm{A}'}$, hence $d_{\mathrm{A}'}f_i\otimes 1_{\mathrm{K}'}=0$;
this proves that the $d_{\mathrm{A}}f_i\otimes 1_\mathrm{K}$ belong to $\mathrm{N}$.
Moreover, their images in $\Upsilon_{\mathrm{K}'/\mathrm{K}}$ are the $d_\mathrm{K}a_i\otimes 1_{\mathrm{K}'}$, which form a \emph{base} of $\Upsilon_{\mathrm{K}'/\mathrm{K}}$ over $\mathrm{K}'$ \sref{0.21.4.7}; taking into account \sref{0.22.4.4.7}, this finishes proving \sref{0.22.4.4}.
\end{proof}

\end{proof}

\begin{env}[22.4.5]
\label{0.22.4.5}
Now consider two rings $\mathrm{A}$, $\mathrm{B}$ of \emph{characteristic} $p>0$, and two homomorphisms
\[
  \mathrm{A}\xrightarrow{i}\mathrm{B}\xrightarrow{j}\mathrm{A}
\]
\oldpage[0\textsubscript{IV-1}]{197}
satisfying the conditions of \sref{0.21.3.1}.
Let furthermore $\mathfrak{m}$ be an ideal \emph{of square zero} in $\mathrm{A}$; set
\[
  \mathrm{K}=\mathrm{A}/\mathfrak{m},\qquad \mathrm{B}_{(\mathrm{K})}=\mathrm{B}\otimes_\mathrm{A}\mathrm{K}=\mathrm{B}/\mathrm{B}i(\mathfrak{m})
\]
and let
\[
  \varphi:\mathrm{A}\to\mathrm{K},\qquad \psi:\mathrm{B}\to\mathrm{B}_{(\mathrm{K})}
\]
be the canonical homomorphisms.
Since $\varphi$ is surjective, $\Omega^1_{\mathrm{B}_{(\mathrm{K})}/\mathrm{A}}$ identifies canonically with $\Omega^1_{\mathrm{B}_{(\mathrm{K})}/\mathrm{K}}$ \sref{0.20.5.15}.
We note that, since $p\geq 2$, for every $x\in\mathfrak{m}$, one has $j(i(x))=x^p=0$, in other words $j(i(\mathfrak{m}))=0$; by passing to quotients, one deduces from $j$ a homomorphism
\[
  j':\mathrm{B}_{(\mathrm{K})}\to\mathrm{A}
\]
such that, if $i'=\psi\circ i:\mathrm{A}\to\mathrm{B}_{(\mathrm{K})}$, we have
\[
  j'(i'(x))=x^p \quad\text{and}\quad i'(j'(y))=y^p
  \quad\text{for}\quad x\in\mathrm{A}\quad\text{and}\quad y\in\mathrm{B}_{(\mathrm{K})}.
\]
Setting, following the notations of \sref{0.21.3.2}
\[
  \Theta_{\mathrm{B}_{(\mathrm{K})}/\mathrm{A}}=\Omega^1_{\mathrm{B}_{(\mathrm{K})}/\mathrm{A}}\otimes_{\mathrm{B}_{(\mathrm{K})}}\mathrm{A}_{[j']}=\Omega^1_{\mathrm{B}_{(\mathrm{K})}/\mathrm{K}}\otimes_{\mathrm{B}_{(\mathrm{K})}}\mathrm{A}_{[j']}
\]
one has therefore \sref{0.21.3.2.2} a canonical homomorphism
\[
\label{0.22.4.5.1}
  \pi_{\mathrm{B}_{(\mathrm{K})}/\mathrm{A}}:\Theta_{\mathrm{B}_{(\mathrm{K})}/\mathrm{A}}\to\Omega^1_\mathrm{A}.
  \tag{22.4.5.1}
\]

On the other hand, one deduces from $i$ and $j$, by passing to quotients (taking into account that $j(i(\mathfrak{m}))=0$) homomorphisms
\[
  \mathrm{K}\xrightarrow{i_0}\mathrm{B}_{(\mathrm{K})}\xrightarrow{j_0}\mathrm{K}
\]
so that one has the commutative diagram
\[
  \begin{tikzcd}
    \mathrm{A} \ar[r,"i"] & \mathrm{B} \ar[r,"j"] & \mathrm{A}\\
    & \psi\downarrow & \nearrow j' &\\
    \mathrm{K} \ar[r,"i_0"] & \mathrm{B}_{(\mathrm{K})} \ar[r,"j_0"] & \mathrm{K}
  \end{tikzcd}
\]
It is clear furthermore that
\[
  j_0(i_0(s))=s^p \quad\text{and}\quad i_0(j_0(t))=t^p
  \quad\text{for}\quad s\in\mathrm{K}\quad\text{and}\quad t\in\mathrm{B}_{(\mathrm{K})}.
\]

One therefore defines in the same way
\[
  \Theta_{\mathrm{B}_{(\mathrm{K})}/\mathrm{K}}=\Omega^1_{\mathrm{B}_{(\mathrm{K})}/\mathrm{K}}\otimes_{\mathrm{B}_{(\mathrm{K})}}\mathrm{K}_{[j_0]}=\Theta_{\mathrm{B}_{(\mathrm{K})}/\mathrm{A}}\otimes_\mathrm{A}\mathrm{K}
\]
hence, tensorizing \sref{0.22.4.5.1} by $\mathrm{K}$, one obtains a \emph{canonical} $\mathrm{K}$-homomorphism
\[
\label{0.22.4.5.2}
  \pi'_{\mathrm{B}_{(\mathrm{K})}/\mathrm{K}}=\pi_{\mathrm{B}_{(\mathrm{K})}/\mathrm{A}}\otimes 1_\mathrm{K}:\Theta_{\mathrm{B}_{(\mathrm{K})}/\mathrm{K}}\to\Omega^1_\mathrm{A}\otimes_\mathrm{A}\mathrm{K}.
  \tag{22.4.5.2}
\]

\oldpage[0\textsubscript{IV-1}]{198}
It follows immediately from the definitions that the diagram
\[
  \begin{tikzcd}
    \Theta_{\mathrm{B}_{(\mathrm{K})}/\mathrm{K}} \ar[r,"\pi'_{\mathrm{B}_{(\mathrm{K})}/\mathrm{K}}"] \ar[dr,"\pi_{\mathrm{B}_{(\mathrm{K})}/\mathrm{K}}",swap] & \Omega^1_\mathrm{A}\otimes_\mathrm{A}\mathrm{K} \ar[d,"\varphi_{\mathrm{K}/\mathrm{A}}"]\\
    & \Omega^1_\mathrm{K}
  \end{tikzcd}
\]
is commutative, $\pi_{\mathrm{B}_{(\mathrm{K})}/\mathrm{K}}$ being the canonical homomorphism corresponding to $i_0$ and $j_0$ \sref{0.21.3.2.2}.
By \emph{restriction} to $\Xi_{\mathrm{B}_{(\mathrm{K})}/\mathrm{K}}=\Ker(\pi_{\mathrm{B}_{(\mathrm{K})}/\mathrm{K}})$, the homomorphism $\pi'_{\mathrm{B}_{(\mathrm{K})}/\mathrm{K}}$ therefore defines a canonical homomorphism
\[
\label{0.22.4.5.3}
  \Xi_{\mathrm{B}_{(\mathrm{K})}/\mathrm{K}}\to\Ker(\Omega^1_\mathrm{A}\otimes_\mathrm{A}\mathrm{K}\to\Omega^1_\mathrm{K})=\Upsilon_{\mathrm{K}/\mathrm{A}}.
  \tag{22.4.5.3}
\]
\end{env}

\begin{env}[22.4.6]
\label{0.22.4.6}
With the hypotheses and notations of \sref{0.22.4.5}, suppose in addition that $\mathfrak{m}$ be a \emph{maximal} ideal of square zero, hence $\mathrm{K}$ a \emph{field}, and let $\mathrm{V}$ denote the ideal $\mathfrak{m}$ considered as a $\mathrm{K}$-vector space.
As $\mathrm{K}$ is an algebra over its prime field $\mathbf{F}_p$ (so that $\mathrm{A}$ is an algebra over $\mathbf{F}_p$, in the sense that $\Omega^1_\mathrm{A}=\Omega^1_{\mathrm{A}/\mathbf{F}_p}$), one deduces from \sref{0.20.5.14} that one has an \emph{exact sequence}
\[
\label{0.22.4.6.1}
  0\to\mathrm{V}\to\Omega^1_\mathrm{A}\otimes_\mathrm{A}\mathrm{K}\to\Omega^1_\mathrm{K}\to 0.
  \tag{22.4.6.1}
\]

One has therefore by \sref{0.22.4.5.3} a homomorphism denoted
\[
\label{0.22.4.6.2}
  \chi'_{\mathrm{B}/\mathrm{A}}:\Xi_{\mathrm{B}_{(\mathrm{K})}/\mathrm{K}}\to\mathrm{V}
  \tag{22.4.6.2}
\]
that we shall call the \emph{characteristic homomorphism} of the $\mathrm{A}$-algebra $\mathrm{B}$ (relative to $i$, $j$ and $\mathfrak{m}$).
It follows from the definitions that for $\chi'_{\mathrm{B}/\mathrm{A}}$ to be \emph{injective}, it is necessary and sufficient that $\pi'_{\mathrm{B}_{(\mathrm{K})}/\mathrm{K}}$ is injective, the kernel of the latter being obviously contained in $\Xi_{\mathrm{B}_{(\mathrm{K})}/\mathrm{K}}$.
\end{env}

\begin{proposition}[22.4.7]
\label{0.22.4.7}
Let $\mathrm{A}$ be a local Artinian ring whose maximal ideal $\mathfrak{m}$ is of square zero, $\mathrm{K}=\mathrm{A}/\mathfrak{m}$ its residue field, $\mathrm{V}$ the ideal $\mathfrak{m}$ considered as a $\mathrm{K}$-vector space; we suppose that $\mathrm{A}$ is of characteristic $p>0$; let
\[
\label{0.22.4.7.1}
  \mathrm{F}_i(\mathrm{T}_i)=\mathrm{T}_i^p-f_i\qquad\text{with}\quad f_i\in\mathrm{A}\quad(1\leq i\leq r)
  \tag{22.4.7.1}
\]
and let $\mathrm{B}$ be the quotient ring of $\mathrm{A}[\mathrm{T}_1,\ldots,\mathrm{T}_r]$ by the ideal generated by the $r$ polynomials $\mathrm{F}_i$.
Then:
\begin{enumerate}
  \item[\rm{(i)}] $\mathrm{B}$ is a local ring.
  \item[\rm{(ii)}] If $\mathfrak{m}'$ is the maximal ideal of $\mathrm{B}$, $\mathrm{K}'=\mathrm{B}/\mathfrak{m}'$ its residue field, $\mathrm{V}'$ the $\mathrm{B}$-module $\mathfrak{m}'/\mathfrak{m}'^2$ considered as a $\mathrm{K}'$-vector space, then, for $\rg_\mathrm{K}(\mathrm{V})=\rg_{\mathrm{K}'}(\mathrm{V}')$, it is necessary and sufficient that the characteristic homomorphism of $\mathrm{K}$-vector spaces
\[
\label{0.22.4.7.2}
  \chi'_{\mathrm{B}/\mathrm{A}}:\Xi_{\mathrm{B}_{(\mathrm{K})}/\mathrm{K}}\to\mathrm{V}
  \tag{22.4.7.2}
\]
  (cf.~\sref{0.22.4.6.2}) is injective.
\end{enumerate}
\end{proposition}

\begin{proof}
\oldpage[0\textsubscript{IV-1}]{199}
Let us again denote by $a_i$ $(1\leq i\leq r)$ the class of $f_i$ mod.~$\mathfrak{m}$; the $a_i$ are not necessarily $p$-independent over $\mathrm{K}^p$, but one can always suppose that the subfamily $(a_i)_{s+1\leq i\leq r}$ is $p$-free over $\mathrm{K}^p$ and forms a $p$-base of $\mathrm{K}^p$ of the field $\mathrm{K}^p(a_1,\ldots,a_r)$ so that
\[
  a_i\in\mathrm{K}^p(a_{s+1},\ldots,a_r)\qquad\text{for}\quad 1\leq i\leq s.
\]

Let us denote by $\mathrm{A}''$ the quotient ring of $\mathrm{A}[\mathrm{T}_{s+1},\ldots,\mathrm{T}_r]$ by the ideal generated by the $\mathrm{F}_i$ of index $i\geq s+1$; then, by \sref{0.22.4.4}, $\mathrm{A}''$ is a local ring whose maximal ideal is $\mathfrak{m}''=\mathfrak{m}\mathrm{A}''$ and the residue field is $\mathrm{K}''=\mathrm{K}(a_{s+1}^{1/p},\ldots,a_r^{1/p})$.
One has thus $a_i\in\mathrm{K}''^p$ for $1\leq i\leq s$ and hence, in $\mathrm{A}''$, the element $f_i$, for $1\leq i\leq s$, can be written
\[
\label{0.22.4.7.3}
  f_i=g_i^p+h_i
  \tag{22.4.7.3}
\]
where $g_i\in\mathrm{A}''$ and $h_i\in\mathfrak{m}''$ (since $\mathfrak{m}''^2=0$ and $p\mathfrak{m}''=0$, it is immediate that $h_i$ is uniquely determined by these conditions).
Replacing $\mathrm{T}_i$ by $\mathrm{T}_i+g_i$, one sees thus that if one sets
\[
  \mathrm{G}_i(\mathrm{T}_i)=\mathrm{T}_i^p+\sum_{k=1}^{p-1}\binom{p}{k}g_i^{p-k}\mathrm{T}_i^k-h_i\qquad(1\leq i\leq s)
\]
$\mathrm{B}$ is the quotient ring of $\mathrm{A}''[\mathrm{T}_1,\ldots,\mathrm{T}_s]$ by the ideal generated by the $\mathrm{G}_i$; but, all the coefficients of $\mathrm{G}_i$ except the leading coefficient belong to $\mathfrak{m}''$, we are in the situation of \sref{0.22.4.3}, which proves (i) (which indeed follows directly from the fact that $\mathrm{B}^p\subset\mathrm{A}$, and by consequence $\mathrm{B}$ has only a single ideal above $\mathfrak{m}$), and furthermore the residue field $\mathrm{K}'$ is \emph{isomorphic} to $\mathrm{K}''$.

\oldpage[0\textsubscript{IV-1}]{200}
Now let us note that for \sref{0.22.4.7.2} to be injective, it is necessary and sufficient that
\[
  \pi'_{\mathrm{B}_{(\mathrm{K})}/\mathrm{K}}:\Omega^1_{\mathrm{B}_{(\mathrm{K})}/\mathrm{K}}\otimes_{\mathrm{B}_{(\mathrm{K})}}\mathrm{K}\to\Omega^1_\mathrm{A}\otimes_\mathrm{A}\mathrm{K}
\]
is injective (by \sref{0.22.4.6}).
But $\mathrm{B}_{(\mathrm{K})}=\mathrm{B}/\mathfrak{m}\mathrm{B}$ is the quotient algebra $\mathrm{C}=\mathrm{K}[\mathrm{T}_1,\ldots,\mathrm{T}_r]$ by the ideal generated by the polynomials $\overline{\mathrm{F}}_i=\mathrm{T}_i^p-a_i$, and since $d_{\mathrm{C}/\mathrm{K}}\overline{\mathrm{F}}_i=0$ it follows from \sref{0.20.5.13} that $\Omega^1_{\mathrm{B}_{(\mathrm{K})}/\mathrm{K}}$ is a free $\mathrm{B}_{(\mathrm{K})}$-module with base $d_{\mathrm{B}_{(\mathrm{K})}/\mathrm{K}}t_i$, where $t_i$ is the canonical image of $\mathrm{T}_i$.
But by definition \sref{0.22.4.5} the canonical image by $j':\mathrm{B}_{(\mathrm{K})}\to\mathrm{A}$ of an element $x\in\mathrm{B}$ whose class mod.~$\mathfrak{m}\mathrm{B}$ is the element $x^p\in\mathrm{A}$; one deduces immediately, since $t_i^p=f_i$, that the image of $d_{\mathrm{B}_{(\mathrm{K})}/\mathrm{K}}t_i\otimes 1$ by $\pi'_{\mathrm{B}_{(\mathrm{K})}/\mathrm{K}}$ is $d_\mathrm{A}f_i\otimes 1$ in $\Omega^1_\mathrm{A}\otimes_\mathrm{A}\mathrm{K}$; the condition that $\pi'_{\mathrm{B}_{(\mathrm{K})}/\mathrm{K}}$ is injective is thus equivalent to the fact that the $d_\mathrm{A}f_i\otimes 1_\mathrm{K}$ are \emph{linearly independent over $\mathrm{K}$}, or equivalently that their images $d_\mathrm{A}f_i\otimes 1_{\mathrm{K}'}$ are \emph{linearly independent over $\mathrm{K}'$} in $\Omega^1_\mathrm{A}\otimes_\mathrm{A}\mathrm{K}'$.

Now, if one applies \sref{0.22.4.4} to the $\mathrm{A}$-algebra $\mathrm{A}''$, one sees that the kernel of the canonical homomorphism $\Omega^1_\mathrm{A}\otimes_\mathrm{A}\mathrm{K}'\to\Omega^1_{\mathrm{A}''}\otimes_{\mathrm{A}''}\mathrm{K}'$ has as its \emph{basis} the $d_\mathrm{A}f_i\otimes 1_{\mathrm{K}'}$ for $s+1\leq i\leq r$.
The preceding condition thus also amounts to saying that the $d_{\mathrm{A}''}h_i$ for $1\leq i\leq s$ are linearly independent in $\Omega^1_{\mathrm{A}''}\otimes_{\mathrm{A}''}\mathrm{K}'$ (taking into account that $d_{\mathrm{A}''}f_i=d_{\mathrm{A}''}h_i$ by \sref{0.22.4.7.3});
on the other hand, by \sref{0.22.4.4}, (ii), if $\mathrm{V}''$ is the $\mathrm{K}'$-vector space $\mathfrak{m}''/\mathfrak{m}''^2$, we have $\rg_{\mathrm{K}'}(\mathrm{V}'')=\rg_\mathrm{K}(\mathrm{V})$ \sref{0.22.4.1}, (ii), and the condition $\rg_{\mathrm{K}'}(\mathrm{V}')=\rg_\mathrm{K}(\mathrm{V})$ is thus \emph{equivalent} to $\rg_{\mathrm{K}'}(\mathrm{V}')=\rg_{\mathrm{K}'}(\mathrm{V}'')$.
But it follows from \sref{0.22.4.3} that this last relation is \emph{equivalent} to the fact that the classes $h_i$
\oldpage[0\textsubscript{IV-1}]{201}
are linearly independent over $\mathrm{K}'$ in $\mathrm{V}''$.
But, in $\Omega^1_{\mathrm{A}''}\otimes_{\mathrm{A}''}\mathrm{K}'$, the images of the $h_i$ by the injection $\mathrm{V}''\to\Omega^1_{\mathrm{A}''}\otimes_{\mathrm{A}''}\mathrm{K}'$ \sref{0.20.5.11.2} are the $d_{\mathrm{A}''}h_i\otimes 1_{\mathrm{K}'}$ (cf.~\sref{0.22.4.6.1}), and this finishes the proof of \sref{0.22.4.7}.
\end{proof}

\begin{corollary}[22.4.7.4]
\label{0.22.4.7.4}
For $\rg_\mathrm{K}(\mathrm{V})=\rg_{\mathrm{K}'}(\mathrm{V}')$, it is necessary and sufficient that the elements $d_\mathrm{A}f_i\otimes 1_\mathrm{K}$ $(1\leq i\leq r)$ be linearly independent in the $\mathrm{K}$-vector space $\Omega^1_\mathrm{A}\otimes_\mathrm{A}\mathrm{K}$.
\end{corollary}

\begin{remark}[22.4.8]
\label{0.22.4.8}
One can define the homomorphisms \sref{0.22.4.5.2} and \sref{0.22.4.6.2} in slightly different conditions.
Let $\mathrm{A}$, $\mathrm{B}$ be two rings, $\mathfrak{m}$ an ideal \emph{of square zero} in $\mathrm{A}$, $\mathrm{K}=\mathrm{A}/\mathfrak{m}$, $p$ a prime number such that $p.1\in\mathfrak{m}$ (in other words, $\mathrm{K}$ is of characteristic $p$, but not necessarily $\mathrm{A}$).
Let furthermore $\mathrm{B}$ be an $\mathrm{A}$-algebra that is an $\mathrm{A}$-module \emph{faithfully flat} (so that $\mathrm{A}\subset\mathrm{B}$), and suppose that we have
\[
\label{0.22.4.8.1}
  y^p\in\mathrm{A}+p\mathrm{B}\subset\mathrm{A}+\mathfrak{m}\mathrm{B}
  \tag{22.4.8.1}
\]
for every $y\in\mathrm{B}$.
If $y^p=x+pz$ with $x\in\mathrm{A}$, $z\in\mathrm{B}$, the \emph{class}$(A\cap pB)$ of $x$ is well determined, and since $\mathrm{A}\cap p\mathrm{B}=p\mathrm{A}$ as $\mathrm{B}$ is a faithfully flat $\mathrm{A}$-module, one has thus defined a canonical map $j:\mathrm{B}\to\mathrm{A}/p\mathrm{A}$, such that, for $x\in\mathrm{A}$, $j(x)$ is the class of $x^p$ mod.~$p\mathrm{A}$.
Furthermore, if $y'\in\mathrm{B}$ and $y'^p=x'+pz'$ with $x'\in\mathrm{A}$, $z'\in\mathrm{B}$, it is immediate that the elements $(y+y')^p-x-x'$ and $(yy')^p-xx'$ belong to $p\mathrm{B}$ by virtue of the fact that the binomial coefficients $\binom{p}{i}$ are multiples of $p$ for $1\leq i\leq p-1$.
The map $j$ is therefore a \emph{ring homomorphism}.
By composition, one deduces from $j$ a homomorphism $j':\mathrm{B}\xrightarrow{j}\mathrm{A}/p\mathrm{A}\to\mathrm{A}/\mathfrak{m}=\mathrm{K}$ and hence a homomorphism $j_0:\mathrm{B}_{(\mathrm{K})}\to\mathrm{K}$; moreover, if $i_0:\mathrm{K}\to\mathrm{B}_{(\mathrm{K})}$ is the canonical injection, $i_0$ and $j_0$ satisfy the conditions of \sref{0.21.3.1} for the rings \emph{of characteristic} $p$, $\mathrm{K}$ and $\mathrm{B}_{(\mathrm{K})}$.
This being so, in order to define a $\mathrm{K}$-homomorphism $\pi'_{\mathrm{B}_{(\mathrm{K})}/\mathrm{K}}:\Omega^1_{\mathrm{B}_{(\mathrm{K})}/\mathrm{K}}\otimes_{\mathrm{B}_{(\mathrm{K})}}\mathrm{K}_{[j_0]}\to\Omega^1_\mathrm{A}\otimes_\mathrm{A}\mathrm{K}$, it suffices \sref{0.20.4.8} to define a $\mathrm{K}$-\emph{derivation}
\[
\label{0.22.4.8.2}
  \mathrm{D}_0:\mathrm{B}_{(\mathrm{K})}\to\Omega^1_\mathrm{A}\otimes_\mathrm{A}\mathrm{K}
  \tag{22.4.8.2}
\]
where the second member is considered as a $\mathrm{B}_{(\mathrm{K})}$-module by means of $j_0$.
For this, it suffices to define an $\mathrm{A}$-\emph{derivation}
\[
\label{0.22.4.8.3}
  \mathrm{D}:\mathrm{B}\to\Omega^1_\mathrm{A}\otimes_\mathrm{A}\mathrm{K}
  \tag{22.4.8.3}
\]
that \emph{vanishes on} $\mathfrak{m}\mathrm{B}$ (the second member being considered as a $\mathrm{B}$-module by means of $j'$).
Indeed, if $y^p=x+pz$ with $y\in\mathrm{B}$, $x\in\mathrm{A}$, $z\in\mathrm{B}$, the element $d_\mathrm{A}x\otimes 1_\mathrm{K}$ of $\Omega^1_\mathrm{A}\otimes_\mathrm{A}\mathrm{K}$ depends only on the \emph{class} $j(y)$ mod.~$p\mathrm{A}$ of $x$ (for since $d_\mathrm{A}(pt)=pd_\mathrm{A}(t)$, hence $d_\mathrm{A}(pt)\otimes 1=0$ in $\Omega^1_\mathrm{A}\otimes_\mathrm{A}\mathrm{K}=\Omega^1_\mathrm{A}/\mathfrak{m}\Omega^1_\mathrm{A}$, since $p.1\in\mathfrak{m}$).
One can therefore take
\[
  \mathrm{D}(y)=d_\mathrm{A}x\otimes 1_\mathrm{K}.
\]

The fact that $\mathrm{D}$ is a derivation follows easily from the fact that $j$ is a ring homomorphism; moreover, for $t\in\mathrm{A}$, one has $d_\mathrm{A}(t^px)=t^pd_\mathrm{A}(x)+pt^{p-1}xd_\mathrm{A}(t)$ and one deduces that $\mathrm{D}(ty)=t\mathrm{D}(y)$, hence $\mathrm{D}$ is an $\mathrm{A}$-derivation; finally, if in addition $t\in\mathfrak{m}$, $\mathrm{D}(ty)=t\mathrm{D}(y)=0$ since $\Omega^1_\mathrm{A}\otimes_\mathrm{A}\mathrm{K}$ is annihilated by $\mathfrak{m}$.
We have thus defined the desired homomorphism \sref{0.22.4.5.2}, and it is immediate to verify that the composition
\[
  \Theta_{\mathrm{B}_{(\mathrm{K})}/\mathrm{K}}\xrightarrow{\pi'_{\mathrm{B}_{(\mathrm{K})}/\mathrm{K}}}\Omega^1_\mathrm{A}\otimes_\mathrm{A}\mathrm{K}\to\Omega^1_\mathrm{K}
\]
is again the canonical homomorphism $\pi_{\mathrm{B}_{(\mathrm{K})}/\mathrm{K}}$ analog to \sref{0.22.4.5.3}.
The analog of \sref{0.22.4.6.2} is also defined in the same way.

Now if $\mathrm{K}$ is a \emph{field}, $\mathrm{A}/p\mathrm{A}$ is an $\mathbf{F}_p$-algebra, and one has the exact sequence \sref{0.20.5.12.1}
\[
  p\mathrm{A}\to\Omega^1_\mathrm{A}\otimes_\mathrm{A}(\mathrm{A}/p\mathrm{A})\to\Omega^1_{\mathrm{A}/p\mathrm{A}}\to 0
\]
which shows, tensorizing with $\mathrm{K}$, that $\Omega^1_\mathrm{A}\otimes_\mathrm{A}\mathrm{K}$ and $\Omega^1_{\mathrm{A}/p\mathrm{A}}\otimes_{\mathrm{A}/p\mathrm{A}}\mathrm{K}$ are isomorphic.
The analog of \sref{0.22.4.6.2} is therefore a homomorphism
\[
\label{0.22.4.8.4}
  \Xi_{\mathrm{B}_{(\mathrm{K})}/\mathrm{K}}\to\mathrm{V}/p\mathrm{A}=\mathfrak{m}/(\mathfrak{m}^2+p\mathrm{A}).
  \tag{22.4.8.4}
\]

That said, the \emph{criterion} \sref{0.22.4.7}, in which one replaces $\mathrm{V}$ and $\mathrm{V}'$ by $\mathrm{V}/p\mathrm{A}$ and $\mathrm{V}'/p\mathrm{B}$ respectively, \emph{remains valid assuming only that $\mathrm{K}$ is of characteristic $p>0$} (in other words, that $p.1\in\mathfrak{m}$ in $\mathrm{A}$), \emph{and that the $\mathrm{F}_i(\mathrm{T}_i)$ are of the more general form} \sref{0.22.4.4.1}: indeed, $\mathrm{B}$ is a free $\mathrm{A}$-module (hence faithfully flat), and it suffices to redo the argument of \sref{0.22.4.7}, replacing \sref{0.22.4.6.2} by \sref{0.22.4.8.4} and $\mathrm{V}''$ by $\mathrm{V}''/p\mathrm{A}''$.
\end{remark}

\subsection{Geometrically regular algebras and formally smooth algebras}
\label{subsection:0.22.5}

\oldpage[0\textsubscript{IV-1}]{201}
\begin{proposition}[22.5.1]
\label{0.22.5.1}
Let $\mathrm{A}$ be a regular local ring, $\mathrm{A}'$ a local ring containing $\mathrm{A}$ and being a finite $\mathrm{A}$-algebra, $\mathfrak{m}$ (resp.~$\mathfrak{m}'$) the maximal ideal of $\mathrm{A}$ (resp.~$\mathrm{A}'$), $\mathrm{K}=\mathrm{A}/\mathfrak{m}$ (resp.~$\mathrm{K}'=\mathrm{A}'/\mathfrak{m}'$) its residue field.
For $\mathrm{A}'$ to be a regular ring, it is necessary and sufficient that
\[
\label{0.22.5.1.1}
  \rg_\mathrm{K}(\mathfrak{m}/\mathfrak{m}^2)=\rg_{\mathrm{K}'}(\mathfrak{m}'/\mathfrak{m}'^2).
  \tag{22.5.1.1}
\]
\end{proposition}

\begin{proof}
Indeed, the first member of \sref{0.22.5.1.1} is $\dim(\mathrm{A})$ \sref{0.17.1.1} and since $\mathrm{A}'\supset\mathrm{A}$ and $\mathrm{A}'$ is a finite $\mathrm{A}$-algebra, $\dim(\mathrm{A}')=\dim(\mathrm{A})$ \sref{0.16.1.5}; the equality \sref{0.22.5.1.1} is therefore necessary and sufficient for $\mathrm{A}'$ to be regular \sref{0.17.1.1}.
\end{proof}

\begin{remark}[22.5.2]
\label{0.22.5.2}
\begin{enumerate}
  \item[\rm{(i)}] Let $\mathrm{A}_1=\mathrm{A}/\mathfrak{m}^2$, $\mathrm{A}_1'=\mathrm{A}'\otimes_\mathrm{A}\mathrm{A}_1=\mathrm{A}'/\mathfrak{m}^2\mathrm{A}'$; as $\mathfrak{m}^2\mathrm{A}'\subset\mathfrak{m}'^2$, $\mathfrak{m}_1'=\mathfrak{m}'/\mathfrak{m}^2\mathrm{A}'$ is the maximal ideal of $\mathrm{A}_1'$ and $\mathfrak{m}_1'/\mathfrak{m}_1'^2$ is a $\mathrm{K}'$-vector space isomorphic to $\mathfrak{m}'/\mathfrak{m}'^2$; the condition of regularity of $\mathrm{A}'$ depends therefore only on the structure of the $\mathrm{A}_1$-\emph{algebra} $\mathrm{A}_1'$, which will allow us below to apply the preliminary results of \sref{0.22.4} on algebras finite over Artinian rings.
  \item[\rm{(ii)}] It follows from the proof of \sref{0.22.5.1} and from \sref{0.16.2.6.1} that one has in all cases
\[
\label{0.22.5.2.1}
  \rg_\mathrm{K}(\mathfrak{m}/\mathfrak{m}^2)\leq\rg_{\mathrm{K}'}(\mathfrak{m}'/\mathfrak{m}'^2).
  \tag{22.5.2.1}
\]
  \item[\rm{(iii)}] Suppose that $\mathrm{A}$ and $\mathrm{A}'$ verify the general hypotheses of \sref{0.22.4.1}.
  It is known \sref{0.19.8.10} that there exists a regular local ring $\mathrm{B}$, with maximal ideal $\mathfrak{n}$, such that $\mathrm{A}$ is isomorphic to $\mathrm{B}/\mathfrak{n}^2$; let us denote by $\mathrm{G}_i\in\mathrm{B}[\mathrm{T}_i]$ a unitary polynomial whose canonical image in $\mathrm{A}[\mathrm{T}_i]$ is $\mathrm{F}_i$ $(1\leq i\leq r)$; then, if $\mathrm{B}'$ is the quotient ring of $\mathrm{B}[\mathrm{T}_1,\ldots,\mathrm{T}_r]$ by the ideal generated by the $\mathrm{G}_i$, it is clear that $\mathrm{B}'$ is a $\mathrm{B}$-module of rank $d$ and $\mathrm{A}'=\mathrm{B}'\otimes_\mathrm{B}\mathrm{A}=\mathrm{B}'/\mathfrak{n}^2\mathrm{B}'$; as $\mathrm{A}'$ is assumed to be a local ring, so is $\mathrm{B}'$, and so $\mathrm{B}'$ is isomorphic to $\mathrm{K}$ and $\mathfrak{n}'/\mathfrak{n}'^2$ isomorphic to $\mathfrak{m}'/\mathfrak{m}'^2$ as $\mathrm{K}'$-vector spaces;
  the inequality \sref{0.22.5.2.1} shows thus that we have, under the hypotheses of \sref{0.22.4.1}
\[
\label{0.22.5.2.2}
  \rg_\mathrm{K}(\mathrm{V})\leq\rg_{\mathrm{K}'}(\mathfrak{m}'/\mathfrak{m}'^2).
  \tag{22.5.2.2}
\]
\end{enumerate}
\end{remark}

\oldpage[0\textsubscript{IV-1}]{202}
\begin{corollary}[22.5.3]
\label{0.22.5.3}
Let $\mathrm{A}$ be a regular local ring, $\mathfrak{m}$ its maximal ideal, $\mathrm{T}_i$ $(1\leq i\leq r)$ indeterminates, and for each $i$, let
\[
\label{0.22.5.3.1}
  \mathrm{F}_i(\mathrm{T}_i)=\mathrm{T}_i^{d_i}+\sum_{1\leq k\leq d_i}c_{ik}\mathrm{T}_i^{d_i-k}
  \tag{22.5.3.1}
\]
be a unitary polynomial of $\mathrm{A}[\mathrm{T}_i]$, with $d_i\geq 2$ for every $i$ and $c_{ik}\in\mathfrak{m}$.
Let $\mathrm{A}'$ be the quotient of the polynomial ring $\mathrm{A}[\mathrm{T}_1,\ldots,\mathrm{T}_r]$ by the ideal $\mathfrak{b}$ generated by the $r$ polynomials $\mathrm{F}_i$.
Then:
\begin{enumerate}
  \item[\rm{(i)}] $\mathrm{A}'$ is a local ring, and its residue field $\mathrm{K}'$ is isomorphic to the residue field $\mathrm{K}$ of $\mathrm{A}$.
  \item[\rm{(ii)}] For $\mathrm{A}'$ to be regular, it is necessary and sufficient that the classes $\xi_i$ mod.~$\mathfrak{m}^2$ of the elements $c_{i,d_i}$ $(1\leq i\leq r)$ be linearly independent over $\mathrm{K}$.
\end{enumerate}
\end{corollary}

\begin{proof}
Indeed, with the notations of \sref{0.22.5.2}, (i), if $\overline{\mathrm{F}}_i$ denotes the polynomial of $\mathrm{A}_1[\mathrm{T}_i]$ obtained by applying to the coefficients of $\mathrm{F}_i$ the homomorphism $\mathrm{A}\to\mathrm{A}_1$, $\mathrm{A}_1'$ is defined from $\mathrm{A}_1$ and the polynomials $\overline{\mathrm{F}}_i$ as in \sref{0.22.4.1}, and the conclusion follows from \sref{0.22.5.1} and \sref{0.22.4.3}.
\end{proof}

\begin{theorem}[22.5.4]
\label{0.22.5.4}
Let $\mathrm{A}$ be a regular local ring, $\mathrm{K}$ its residue field; suppose that $\mathrm{K}$ is of characteristic $p>0$; let
\[
\label{0.22.5.4.1}
  \mathrm{F}_i(\mathrm{T}_i)=\mathrm{T}_i^p+p\sum_{k=1}^{p-1}c_{ik}\mathrm{T}_i^{p-k}-f_i
  \tag{22.5.4.1}
\]
with $c_{ik}\in\mathrm{A}$ and $f_i\in\mathrm{A}$ $(1\leq i\leq r)$; let $\mathrm{B}$ be the quotient ring of $\mathrm{A}[\mathrm{T}_1,\ldots,\mathrm{T}_r]$ by the ideal generated by the $r$ polynomials $\mathrm{F}_i$.
Then $\mathrm{B}$ is a local ring, and the following conditions are equivalent:
\begin{enumerate}
  \item[a)] $\mathrm{B}$ is regular;
  \item[b)] the elements $d_\mathrm{A}f_i\otimes 1_\mathrm{K}$ $(1\leq i\leq r)$ are linearly independent in the $\mathrm{K}$-vector space $\Omega^1_\mathrm{A}\otimes_\mathrm{A}\mathrm{K}$;
  \item[c)] if $\mathrm{B}'=\mathrm{B}/\mathfrak{m}^2\mathrm{B}$ and if one sets $\Xi_{\mathrm{B}_{(\mathrm{K})}/\mathrm{K}}=\Xi_{\mathrm{B}'_{(\mathrm{K})}/\mathrm{K}}$ \sref{0.22.4.5}, the characteristic homomorphism $\Xi_{\mathrm{B}_{(\mathrm{K})}/\mathrm{K}}\to\mathfrak{m}/\mathfrak{m}^2$ \sref{0.22.4.6.2} is injective.
\end{enumerate}
\end{theorem}

\begin{proof}
Indeed, $\mathrm{B}'$ is defined again from $\mathrm{A}_1=\mathrm{A}/\mathfrak{m}^2$ and the polynomials $\overline{\mathrm{F}}_i$ obtained by applying the homomorphism $\mathrm{A}\to\mathrm{A}_1$ to the coefficients.
Taking into account that $\Omega^1_\mathrm{A}\otimes_\mathrm{A}\mathrm{K}$ is canonically isomorphic to $\Omega^1_{\mathrm{A}_1}\otimes_{\mathrm{A}_1}\mathrm{K}$ \sref{0.20.5.12}, (ii), the theorem follows from \sref{0.22.5.1} and \sref{0.22.4.7.4}, using remark \sref{0.22.4.8} when $\mathrm{A}_1$ is not of characteristic $p$.
\end{proof}

\begin{env}[22.5.5]
\label{0.22.5.5}
Let $k$ be a field of characteristic $p>0$, $\mathrm{A}$ a local ring, $k'$ a \emph{finite} radiciel extension of $k$ such that $k'^p\subset k$, $\mathfrak{m}$ the maximal ideal of $\mathrm{A}$, $\mathrm{K}=\mathrm{A}/\mathfrak{m}$ its residue field; recall that $\mathrm{A}'=\mathrm{A}\otimes_kk'$ is a local ring whose residue field is the same as that of $\mathrm{K}\otimes_kk'$ (Bourbaki, \emph{Alg.~comm.}, chap.~V, \textsection~2, lemma~4).
\oldpage[0\textsubscript{IV-1}]{203}
We note that one has $\mathrm{A}'_{(\mathrm{K})}=\mathrm{A}'\otimes_\mathrm{A}\mathrm{K}=\mathrm{K}\otimes_kk'$, and hence \sref{0.21.3.6}
\[
\label{0.22.5.5.1}
  \Theta_{\mathrm{A}'_{(\mathrm{K})}/\mathrm{K}}=\Theta_{k'/k}\otimes_k\mathrm{K}
  \tag{22.5.5.1}
\]
to a canonical isomorphism.
Furthermore, it is known \sref{0.21.4.7} that one has $\Xi_{k'/k}=0$, in other words the canonical homomorphism \sref{0.21.3.2.2}
\[
  \pi_{k'/k}:\Theta_{k'/k}\to\Omega^1_k
\]
is injective, so that one has a canonical \emph{injection}
\[
\label{0.22.5.5.2}
  \Theta_{\mathrm{A}'_{(\mathrm{K})}/\mathrm{K}}\to\Omega^1_k\otimes_k\mathrm{K}
  \tag{22.5.5.2}
\]
which can be described explicitly as follows (taking into account \sref{0.22.5.5.1}): for every $x'\in k'$, to the element $d_{k'/k}(x'\otimes 1_\mathrm{K})\otimes 1_\mathrm{K}$ of $\Omega^1_{\mathrm{A}'_{(\mathrm{K})}/\mathrm{K}}\otimes_{\mathrm{A}'_{(\mathrm{K})}}\mathrm{K}$, one associates the element $d_k(x'^p)\otimes 1_\mathrm{K}$ of $\Omega^1_k\otimes_k\mathrm{K}$.
\end{env}

\begin{lemma}[22.5.6]
\label{0.22.5.6}
With the notations of \sref{0.22.5.5}:
\begin{enumerate}
  \item[\rm{(i)}] The kernel of the canonical homomorphism $\pi_{\mathrm{A}'_{(\mathrm{K})}/\mathrm{K}}$ \sref{0.21.3.2.2} identifies by \sref{0.22.5.5.2} with $(\Theta_{k'/k}\otimes_k\mathrm{K})\cap\Upsilon_{\mathrm{K}/k}$.
  \item[\rm{(ii)}] Let $\mathrm{A}_1=\mathrm{A}/\mathfrak{m}^2$, $\mathrm{A}_1'=\mathrm{A}_1\otimes_kk'=\mathrm{A}'/\mathfrak{m}^2\mathrm{A}'$; then the characteristic homomorphism $\chi'_{\mathrm{A}_1'/\mathrm{A}_1}$ \sref{0.22.4.6.2} identifies with the restriction of the characteristic homomorphism $\chi_{\mathrm{A}/k}:\Upsilon_{\mathrm{K}/k}\to\mathrm{V}=\mathfrak{m}/\mathfrak{m}^2$ to $(\Theta_{k'/k}\otimes_k\mathrm{K})\cap\Upsilon_{\mathrm{K}/k}$ \sref{0.20.6.24}.
\end{enumerate}
\end{lemma}

\begin{proof}
(i) By virtue of \sref{0.20.6.9}, applied replacing $\mathrm{A}$ by the prime field, $\mathrm{C}$ by $\mathrm{K}$ and $\mathrm{E}$ by $\mathrm{A}_1=\mathrm{A}/\mathfrak{m}^2$ (extension of $\mathrm{K}$ by $k$, $\mathrm{V}=\mathfrak{m}/\mathfrak{m}^2$), it suffices (cf.~\sref{0.22.5.5}) to verify that if $x'\in k'$ and if $z\in\mathrm{A}_1$ is an element whose class mod.~$\mathfrak{m}/\mathfrak{m}^2$ is $x'^p\in k$, then the image $d_k(x'^p)\otimes 1_\mathrm{K}$ by the canonical homomorphism $\Omega^1_k\otimes_k\mathrm{K}\to\Omega^1_{\mathrm{A}_1}\otimes_{\mathrm{A}_1}\mathrm{K}$ is $d_{\mathrm{A}_1}z\otimes 1_\mathrm{K}$, which follows from the definitions.

(ii) As $\mathrm{A}'_{(\mathrm{K})}=\mathrm{K}\otimes_kk'=(\mathrm{A}_1')_{(\mathrm{K})}$ (by virtue of the relation $\mathrm{K}=\mathrm{A}_1/(\mathfrak{m}/\mathfrak{m}^2)$), $\Theta_{\mathrm{A}'_{(\mathrm{K})}/\mathrm{K}}$ identifies with $\Theta_{(\mathrm{A}_1')_{(\mathrm{K})}/\mathrm{K}}$, and the verification follows from the same calculation as in (i) and from the definition of $\pi'_{\mathrm{A}'_{(\mathrm{K})}/\mathrm{K}}$.
\end{proof}

\begin{proposition}[22.5.7]
\label{0.22.5.7}
Let $k$ be a field of characteristic $p>0$, $\mathrm{A}$ a $k$-algebra that is a regular local ring, $\mathfrak{m}$ its maximal ideal, $\mathrm{K}=\mathrm{A}/\mathfrak{m}$ its residue field; let $\mathrm{V}=\mathfrak{m}/\mathfrak{m}^2$, considered as a $\mathrm{K}$-vector space.
Let $k'$ be a finite extension of $k$ such that $k'^p\subset k$, and $\mathrm{A}'=\mathrm{A}\otimes_kk'$; for $\mathrm{A}'$ to be a regular local ring, it is necessary and sufficient that the restriction of the characteristic homomorphism $\chi_{\mathrm{A}/k}:\Upsilon_{\mathrm{K}/k}\to\mathrm{V}$ \sref{0.20.6.24} to $(\Theta_{k'/k}\otimes_k\mathrm{K})\cap\Upsilon_{\mathrm{K}/k}$ is injective.
\end{proposition}

\begin{proof}
\oldpage[0\textsubscript{IV-1}]{204}
We must prove that the statement is equivalent to \sref{0.22.5.1.1}, by denoting by $\mathfrak{m}'$ the maximal ideal of $\mathrm{A}'$, with $\mathrm{K}'=\mathrm{A}'/\mathfrak{m}'$ its residue field.
With the notations of \sref{0.22.5.2}, (i), one has $\mathrm{A}_1'=\mathrm{A}_1\otimes_kk'$.
Now, let $(a_i)_{1\leq i\leq r}$ be a $p$-base of $k'^p$ over $k$; the elements $d_k(x_\alpha^p)$ form a base of the $k$-vector space $\Omega^1_k$ \sref{0.21.4.5}, hence the elements $d_k(x_\alpha^p)\otimes 1_\mathrm{K}$ form a base of the $\mathrm{K}$-vector space $\Omega^1_k\otimes_k\mathrm{K}$.
One concludes (cf.~\sref{0.22.5.5}) that when $k'$ ranges over the finite subextensions of $k^{1/p}$, the family of subspaces $\Theta_{k'/k}\otimes_k\mathrm{K}$ of $\Omega^1_k\otimes_k\mathrm{K}$ is filtering increasing and has as union $\Omega^1_k\otimes_k\mathrm{K}$.
It follows then from \sref{0.22.5.7} that the condition $b'$) implies that $\chi_{\mathrm{A}/k}$ is injective, which is none other than condition (ii) of \sref{0.22.2.2}.

One can apply the criterion \sref{0.22.4.7}, and the condition is equivalent to the fact that $\chi'_{\mathrm{A}_1'/\mathrm{A}_1}$ is injective; one concludes with the aid of \sref{0.22.5.6}.
\end{proof}

\begin{theorem}[22.5.8]
\label{0.22.5.8}
Let $k$ be a field, $p$ its characteristic exponent, $\mathrm{A}$ a local Noetherian $k$-algebra.
The following conditions are equivalent:
\begin{enumerate}
  \item[a)] $\mathrm{A}$ is a $k$-algebra formally smooth (for its preadic topology).
  \item[b)] $\mathrm{A}$ is geometrically regular over $k$.
  \item[b')] For every finite extension $k'$ of $k$ such that $k'^p\subset k$, $\mathrm{A}'=\mathrm{A}\otimes_kk'$ is a regular ring.
\end{enumerate}
Taking into account \sref{0.19.6.6}, one can restrict oneself to the case $p>1$.
\end{theorem}

\begin{proof}
The fact that $a$) implies $b$) is none other than \sref{0.19.6.5}, and $b$) trivially implies $b'$).
Let us show that $b'$) implies the conditions of \sref{0.22.2.2} are satisfied; it is evident for the first \sref{0.17.1.1} since $b'$) implies first of all that $\mathrm{A}$ is regular.
Next, let $(x_\alpha)_{\alpha\in\mathrm{I}}$ be a $p$-base of $k^{1/p}$ over $k$; we know that the elements $d_k(x_\alpha^p)$ form a base of the $k$-vector space $\Omega^1_k$ \sref{0.21.4.5}, hence the elements $d_k(x_\alpha^p)\otimes 1_\mathrm{K}$ form a base of the $\mathrm{K}$-vector space $\Omega^1_k\otimes_k\mathrm{K}$.
One concludes (cf.~\sref{0.22.5.5}) that when $k'$ ranges over the finite subextensions of $k^{1/p}$, the family of subspaces $\Theta_{k'/k}\otimes_k\mathrm{K}$ is filtering increasing and has as union $\Omega^1_k\otimes_k\mathrm{K}$.
It follows then from \sref{0.22.5.7} that the condition $b'$) implies that $\chi_{\mathrm{A}/k}$ is injective, which is none other than condition (ii) of \sref{0.22.2.2}.
\end{proof}

\begin{corollary}[22.5.9]
\label{0.22.5.9}
Let $k$ be a field, $\mathrm{A}$ a local Noetherian $k$-algebra that is formally smooth.
Then, for every prime ideal $\mathfrak{p}$ of $\mathrm{A}$, the local ring $\mathrm{A}_\mathfrak{p}$ is a $k$-algebra formally smooth (for the $\mathfrak{p}$-preadic topology).
\end{corollary}

\begin{proof}
Indeed, with the notations of \sref{0.22.5.8}, $b'$), it suffices to show that the local ring $\mathrm{A}_\mathfrak{p}\otimes_kk'$ is regular; but since this is a ring of fractions of the local ring $\mathrm{A}\otimes_kk'$ and the latter is regular by hypothesis, it is the same by \sref{0.17.3.6}.
\end{proof}

\oldpage[0\textsubscript{IV-1}]{205}
\begin{remark}[22.5.10]
\label{0.22.5.10}
\begin{enumerate}
  \item[\rm{(i)}] The conclusion of \sref{0.22.5.8} is in default when, in the condition $b'$), one restricts oneself to \emph{monogenic} extensions $k'=k(x)$ of $k$ (with $x^p\in k$).
  One can indeed show that for all $x\in k^{1/p}$ such that $x\notin k$, one has $\Upsilon_{\mathrm{K}/k}\cap(\Theta_{k'/k}\otimes_k\mathrm{K})=0$, which implies that $\Upsilon_{\mathrm{K}/k}\neq 0$; this signifies that for every $x\in k^{1/p}$ such that $x\notin k$, one must have $d_\mathrm{K}(x^p)\otimes 1_\mathrm{K}\notin\Upsilon_{\mathrm{K}/k}$, or yet $d_\mathrm{K}(x^p)\neq 0$, that is to say $x^p\notin\mathrm{K}^p$; in other words, one must have $\mathrm{K}^p\cap k=k^p$, that is, $\mathrm{K}$ is an \emph{inseparable} extension of $k$.
  Now, one easily constructs examples of such extensions: start from a perfect field $k_0$ of characteristic $p>0$, let $\mathrm{X}$, $\mathrm{Y}$, $\mathrm{Z}$ be three indeterminates and set $k=k_0(\mathrm{X},\mathrm{Y})$ and $\mathrm{K}=k(\mathrm{X}^{1/p}+\mathrm{Z}\mathrm{Y}^{1/p})$; one verifies easily that $\mathrm{K}$ satisfies the preceding conditions, hence $\Upsilon_{\mathrm{K}/k}\neq 0$.
  Applying now the lemma \sref{0.22.2.5.2} taking $\mathrm{V}=\mathrm{K}$ and for $h$ the zero homomorphism; $\mathrm{A}$ is then a valuation ring of discrete valuation having $\mathrm{K}$ as residue field, with $\chi_{\mathrm{A}/k}=0$, and is not a $k$-algebra formally smooth by \sref{0.22.2.2}; however $\mathrm{A}\otimes_kk'$ is a regular ring for every \emph{monogenic} subextension $k'\subset k^{1/p}$ of $k$.
  \item[\rm{(ii)}] The corollary \sref{0.22.5.9} leads to the following question: if $\mathrm{A}$ is a Noetherian $k$-algebra, under what conditions is the set of prime ideals $\mathfrak{p}\in\Spec(\mathrm{A})$ such that $\mathrm{A}_\mathfrak{p}$ is a $k$-algebra formally smooth an \emph{open} set?
  We will address this later in certain particular cases.
\end{enumerate}
\end{remark}

\subsection{Zariski's Jacobian criterion}
\label{subsection:0.22.6}

\begin{theorem}[22.6.1]
\label{0.22.6.1}
\textnormal{(Jacobian criterion for formal smoothness.)}
Let $\mathrm{A}$, $\mathrm{B}$ be two topological rings, $u:\mathrm{A}\to\mathrm{B}$ a continuous homomorphism making $\mathrm{B}$ a formally smooth $\mathrm{A}$-algebra, $\mathfrak{J}$ an ideal of $\mathrm{B}$ (not necessarily closed), $\mathrm{C}=\mathrm{B}/\mathfrak{J}$ the quotient topological algebra.
For $\mathrm{C}$ to be a formally smooth $\mathrm{A}$-algebra, it is necessary and sufficient that the canonical homomorphism (cf.~\sref{0.20.5.11.2})
\[
\label{0.22.6.1.1}
  \delta_{\mathrm{C}/\mathrm{B}/\mathrm{A}}:\mathfrak{J}/\mathfrak{J}^2\to\Omega^1_{\mathrm{B}/\mathrm{A}}\otimes_\mathrm{B}\mathrm{C}
  \tag{22.6.1.1}
\]
be \emph{formally invertible on the left} (cf.~\sref{0.19.1.5}).
\end{theorem}

\begin{proof}
Indeed, saying that $\mathrm{B}$ (resp.~$\mathrm{C}$) is a formally smooth $\mathrm{A}$-algebra signifies \sref{0.19.4.4} that $\operatorname{Exalcotop}_\mathrm{A}(\mathrm{B},\mathrm{L})=0$ (resp.~$\operatorname{Exalcotop}_\mathrm{A}(\mathrm{C},\mathrm{L})=0$) for every $\mathrm{B}$-module (resp.~$\mathrm{C}$-module) discrete $\mathrm{L}$ annihilated by an open ideal of $\mathrm{B}$ (resp.~$\mathrm{C}$).
By hypothesis $\operatorname{Exalcotop}_\mathrm{A}(\mathrm{B},\mathrm{L})=0$ for every such discrete $\mathrm{C}$-module $\mathrm{L}$ annihilated by an open ideal of $\mathrm{C}$, saying that $\mathrm{C}$ is a formally smooth $\mathrm{A}$-algebra thus signifies that the canonical homomorphism $\operatorname{Exalcotop}_\mathrm{A}(\mathrm{C},\mathrm{L})\to\operatorname{Exalcotop}_\mathrm{A}(\mathrm{B},\mathrm{L})$ is \emph{injective}, and the theorem follows from \sref{0.20.7.8}.
\end{proof}

Recall \sref{0.19.5.3} that when $\mathrm{B}$ and $\mathrm{C}$ are formally smooth $\mathrm{A}$-algebras, $\mathfrak{J}/\mathfrak{J}^2$ is a $\mathrm{C}$-module topologically \emph{formally projective}; moreover, the canonical homomorphisms $\varphi_n:\mathbf{S}^n_\mathrm{C}(\mathfrak{J}/\mathfrak{J}^2)\to\mathfrak{J}^n/\mathfrak{J}^{n+1}$ are \emph{formal bimorphisms} when $\mathrm{B}$ is assumed to be a preadmissible ring.

\begin{corollary}[22.6.2]
\label{0.22.6.2}
With the hypotheses and notations of \sref{0.22.6.1}, suppose furthermore that the square of every open ideal of $\mathrm{B}$ is open, and that on $\mathfrak{J}$ the induced topology is identical to that deduced from that of $\mathrm{B}$ \sref{0.19.0} (note that these two conditions are satisfied when $\mathrm{B}$ is \emph{Noetherian} and its topology is preadic $(\mathbf{0}_\mathrm{I},~7.3.2)$, or if $\mathrm{B}$ has the $\mathfrak{J}$-preadic topology).
Let $(\mathfrak{b}_\lambda)$ be a fundamental system of open ideals of $\mathrm{B}$ and set $\mathrm{B}_\lambda=\mathrm{B}/\mathfrak{b}_\lambda$ for every $\lambda$.
Then:
\begin{enumerate}
  \item[\rm{(i)}] The following conditions are equivalent:
  \begin{enumerate}
    \item[a)] $\mathrm{C}$ is a formally smooth $\mathrm{A}$-algebra.
    \item[b)] For every $\lambda$, the homomorphism of $\mathrm{B}_\lambda$-modules
\[
\label{0.22.6.2.1}
  (\mathfrak{J}/\mathfrak{J}^2)\otimes_\mathrm{B}\mathrm{B}_\lambda\to\Omega^1_{\mathrm{B}/\mathrm{A}}\otimes_\mathrm{B}\mathrm{B}_\lambda
  \tag{22.6.2.1}
\]
  deduced from $\delta_{\mathrm{C}/\mathrm{B}/\mathrm{A}}$ by tensorization, is invertible on the left (in other words is an isomorphism onto a direct factor of $\Omega^1_{\mathrm{B}/\mathrm{A}}\otimes_\mathrm{B}\mathrm{B}_\lambda$).
  \end{enumerate}
  \item[\rm{(ii)}] Suppose furthermore that the topological ring $\mathrm{C}$ be preadmissible $(\mathbf{0}_\mathrm{I},~7.1.2)$ and
\end{enumerate}
\end{corollary}

\oldpage[0\textsubscript{IV-1}]{206}
let $\mathfrak{L}$ be an ideal of definition of $\mathrm{C}$; then the conditions $a$) and $b$) of (i) are also equivalent to:
\begin{enumerate}
  \item[c)] The homomorphism of $(\mathrm{C}/\mathfrak{L})$-modules
  \[
    (\mathfrak{J}/\mathfrak{J}^2)\otimes_\mathrm{C}(\mathrm{C}/\mathfrak{L})\to\Omega^1_{\mathrm{B}/\mathrm{A}}\otimes_\mathrm{B}(\mathrm{C}/\mathfrak{L})
  \]
  is invertible on the left.
\end{enumerate}

\begin{proof}
The hypotheses imply that on $\Omega^1_{\mathrm{B}/\mathrm{A}}$ the topology is deduced from that of $\mathrm{B}$ \sref{0.20.4.5}, hence (i) follows from \sref{0.22.6.1} and \sref{0.19.1.7}.
On the other hand, recall \sref{0.20.4.9} that $\Omega^1_{\mathrm{B}/\mathrm{A}}$ is a $\mathrm{B}$-module formally \emph{projective}, hence $\Omega^1_{\mathrm{B}/\mathrm{A}}\otimes_\mathrm{B}\mathrm{B}_\lambda$ is a projective $\mathrm{B}_\lambda$-module for every $\lambda$ \sref{0.19.2.4}; the equivalence of $c$) follows then from \sref{0.19.1.9}.
\end{proof}

In particular:
\begin{corollary}[22.6.3]
\label{0.22.6.3}
Let $\mathrm{A}$ be a ring, $\mathrm{B}$ an $\mathrm{A}$-algebra, $\mathfrak{J}$ an ideal of $\mathrm{B}$, $\mathrm{C}=\mathrm{B}/\mathfrak{J}$; we suppose $\mathrm{A}$, $\mathrm{B}$, $\mathrm{C}$ equipped with the discrete topologies and that $\mathrm{B}$ is a formally smooth $\mathrm{A}$-algebra.
For $\mathrm{C}$ to be a formally smooth $\mathrm{A}$-algebra, it is necessary and sufficient that the canonical homomorphism \sref{0.22.6.1.1} be invertible on the left.
\end{corollary}

\begin{proof}
Note that since \emph{every} $\mathrm{A}$-algebra $\mathrm{C}$ is a quotient of a polynomial algebra $\mathrm{A}[\mathrm{T}_\alpha]_{\alpha\in\mathrm{I}}$, and the latter is formally smooth \sref{0.19.3.3}, the criterion \sref{0.22.6.3} allows in principle to determine whether $\mathrm{C}$ is a formally smooth $\mathrm{A}$-algebra.
\end{proof}

\begin{proposition}[22.6.4]
\label{0.22.6.4}
Let $\mathrm{A}$ be a ring, $\mathrm{B}$ a formally smooth $\mathrm{A}$-algebra (for the discrete topologies), $\mathfrak{J}$ an ideal of $\mathrm{B}$, $\mathrm{C}=\mathrm{B}/\mathfrak{J}$; suppose that $\mathfrak{J}/\mathfrak{J}^2$ is a $\mathrm{C}$-module of finite type.
Let $\mathfrak{p}$ be a prime ideal of $\mathrm{C}$, $k(\mathfrak{p})$ the residue field of $\mathrm{C}$ at $\mathfrak{p}$, $\mathfrak{p}'$ the inverse image of $\mathfrak{p}$ in $\mathrm{A}$.
The following conditions are equivalent:
\begin{enumerate}
  \item[a)] $\mathrm{C}_\mathfrak{p}$ is a formally smooth $\mathrm{A}_{\mathfrak{q}}$-algebra (for the discrete topologies).
  \item[a')] $\mathrm{C}_\mathfrak{p}$ is a formally smooth $\mathrm{A}$-algebra (for the discrete topologies).
  \item[b)] The canonical homomorphism
\[
\label{0.22.6.4.1}
  (\mathfrak{J}/\mathfrak{J}^2)\otimes_\mathrm{C}k(\mathfrak{p})\to\Omega^1_{\mathrm{B}/\mathrm{A}}\otimes_\mathrm{B}k(\mathfrak{p})
  \tag{22.6.4.1}
\]
  is injective.
  \item[c)] There exists $f\in\mathrm{C}-\mathfrak{p}$ such that $\mathrm{C}_f$ is a formally smooth $\mathrm{A}$-algebra for the discrete topology.
\end{enumerate}

If furthermore $\mathrm{B}$ is \emph{Noetherian}, the preceding conditions are also equivalent to:
\begin{enumerate}
  \item[d)] $\mathrm{C}_\mathfrak{p}$ is a formally smooth $\mathrm{A}_{\mathfrak{q}}$-algebra for the $\mathfrak{p}$-preadic topology on $\mathrm{C}_\mathfrak{p}$ and the $\mathfrak{q}$-adic topology (or the discrete topology) on $\mathrm{A}_{\mathfrak{q}}$.
\end{enumerate}
\end{proposition}

\begin{proof}
\oldpage[0\textsubscript{IV-1}]{207}
As $\mathrm{A}_{\mathfrak{q}}$ is a formally smooth $\mathrm{A}$-algebra (by \sref{0.19.3.5}, (iv) and (ii)), one already knows that $a$) and $a'$) are equivalent \sref{0.19.3.5}, (ii).
On the other hand $\mathrm{C}_\mathfrak{p}=\mathrm{B}_{\mathfrak{p}'}/\mathfrak{J}\mathrm{B}_{\mathfrak{p}'}$, and $\mathrm{B}_{\mathfrak{p}'}$ is an $\mathrm{A}$-algebra formally smooth for the discrete topology \sref{0.19.3.5}, (iv); moreover $\Omega^1_{\mathrm{B}/\mathrm{A}}$ is a $\mathrm{B}$-module \emph{projective}, hence $\Omega^1_{\mathrm{B}/\mathrm{A}}\otimes_\mathrm{B}\mathrm{B}_\lambda$ is a projective $\mathrm{B}_\lambda$-module for every $\lambda$ \sref{0.19.2.4}; and $\mathfrak{J}/\mathfrak{J}^2$ a $\mathrm{B}$-module of finite type, the equivalence of $a'$) and $b$) follows then from the application of \sref{0.22.6.3} to $\mathrm{B}_{\mathfrak{p}'}$ and $\mathrm{C}_\mathfrak{p}$, taking into account the fact that $\Omega^1_{\mathrm{B}/\mathrm{A}}$ is a projective $\mathrm{B}$-module \sref{0.20.4.9} and $\mathfrak{J}/\mathfrak{J}^2$ a $\mathrm{C}$-module of finite type.
The application of \sref{0.19.1.12} proves also the equivalence of $b$) and $c$).
Finally, let us note that since $\mathrm{B}_{\mathfrak{p}'}$ is a $\mathrm{A}_{\mathfrak{q}}$-algebra formally smooth for the $\mathfrak{p}'$-preadic topology and its topology is preadic \sref{0.19.3.8}, and $\mathfrak{J}/\mathfrak{J}^2$ a $\mathrm{C}$-module of finite type, and $\mathrm{C}_\mathfrak{p}$ preadmissible and $\mathfrak{p}\mathrm{C}_\mathfrak{p}$ is an ideal of definition for its topology, the equivalence of $c$) and $a$) in (i) follows from \sref{0.22.6.4}, and the last assertion follows from \sref{0.22.6.2}.
\end{proof}

\begin{corollary}[22.6.5]
\label{0.22.6.5}
Under the general hypotheses of \sref{0.22.6.4}, the set of $\mathfrak{p}\in\Spec(\mathrm{C})$ such that $\mathrm{C}_\mathfrak{p}$ is a formally smooth $\mathrm{A}_{\mathfrak{q}}$-algebra (or an $\mathrm{A}$-algebra formally smooth), where $\mathfrak{q}$ denotes the inverse image of $\mathfrak{p}$ in $\mathrm{A}$, \emph{for the discrete topologies}, is open in $\Spec(\mathrm{C})$.
\end{corollary}

\begin{proof}
This follows from form $c$) of \sref{0.22.6.4}.
Note that when $\mathrm{B}$ is Noetherian, one can replace the discrete topologies by the preadic topologies.
\end{proof}

\begin{corollary}[22.6.6]
\label{0.22.6.6}
Under the general hypotheses of \sref{0.22.6.4}, the following conditions are equivalent:
\begin{enumerate}
  \item[a)] $\mathrm{C}$ is a formally smooth $\mathrm{A}$-algebra (for the discrete topologies).
  \item[b)] For every $\mathfrak{p}\in\Spec(\mathrm{C})$ (or only for every \emph{maximal} ideal $\mathfrak{p}$ of $\mathrm{C}$), $\mathrm{C}_\mathfrak{p}$ is a formally smooth $\mathrm{A}$-algebra (or a formally smooth $\mathrm{A}_{\mathfrak{q}}$-algebra, where $\mathfrak{q}$ denotes the inverse image of $\mathfrak{p}$ in $\mathrm{A}$) \emph{for the discrete topologies}.
\end{enumerate}
Furthermore, when $\mathrm{B}$ is Noetherian, one can replace in $b$) the discrete topologies by the preadic topologies on $\mathrm{A}_{\mathfrak{q}}$ and $\mathrm{C}_\mathfrak{p}$.
\end{corollary}

\begin{proof}
The last assertion follows from \sref{0.22.6.4}.
By virtue of \sref{0.22.6.3}, the condition $a$) is equivalent to the fact that the homomorphism \sref{0.22.6.1.1} be invertible on the left; the equivalence of $a$) and $b$) follows from the application of \sref{0.22.6.3} to $\mathrm{B}_{\mathfrak{p}'}$ and $\mathrm{C}_\mathfrak{p}$, taking into account the fact that $\Omega^1_{\mathrm{B}/\mathrm{A}}$ is a projective $\mathrm{B}$-module and $\mathfrak{J}/\mathfrak{J}^2$ a $\mathrm{C}$-module of finite type.
\end{proof}

\begin{proposition}[22.6.7]
\label{0.22.6.7}
Let $k_0$ be a field, $k$ a separable extension of $k_0$, $\mathrm{C}$ a $k$-algebra of finite type.
\begin{enumerate}
  \item[\rm{(i)}] Let $\mathfrak{p}$ be a prime ideal of $\mathrm{C}$.
  The following conditions are equivalent:
  \begin{enumerate}
    \item[a)] $\mathrm{C}_\mathfrak{p}$ is a $k_0$-algebra formally smooth for the discrete topology.
    \item[b)] $\mathrm{C}_\mathfrak{p}$ is a $k_0$-algebra formally smooth for the $\mathfrak{p}$-preadic topology.
    \item[c)] There exists $f\in\mathrm{C}-\mathfrak{p}$ such that $\mathrm{C}_f$ is a $k_0$-algebra formally smooth for the discrete topology.
    \item[d)] $\mathrm{C}_\mathfrak{p}$ is a geometrically regular ring over $k_0$ \sref{0.19.6.5}.
  \end{enumerate}
  Furthermore the set of $\mathfrak{p}\in\Spec(\mathrm{C})$ having any of these properties is open in $\Spec(\mathrm{C})$.
  \item[\rm{(ii)}] The following conditions are equivalent:
  \begin{enumerate}
    \item[a)] $\mathrm{C}$ is a $k_0$-algebra formally smooth for the discrete topology.
    \item[b)] Every $\mathfrak{p}\in\Spec(\mathrm{C})$ satisfies the equivalent conditions of (i).
    \item[c)] Every maximal ideal $\mathfrak{m}$ of $\mathrm{C}$ satisfies the equivalent conditions of (i).
    \item[d)] For every extension $k'$ of $k_0$, $\mathrm{C}\otimes_{k_0}k'$ is a regular ring \sref{0.17.3.6}; we also say that $\mathrm{C}$ is a geometrically regular ring over $k_0$.
  \end{enumerate}
\end{enumerate}
\end{proposition}

\begin{proof}
\oldpage[0\textsubscript{IV-1}]{208}
One knows that $\mathrm{C}$ is always isomorphic to a $k$-algebra of the form $\mathrm{B}/\mathfrak{J}$.
Since $\mathrm{B}$ is a $k$-algebra formally smooth \sref{0.19.3.3} and $k$, being separable over $k_0$, is a $k_0$-algebra formally smooth \sref{0.19.6.1}, $\mathrm{B}$ is a $k_0$-algebra formally smooth \sref{0.19.3.5}, (ii).
The equivalence of conditions $a$), $b$), $c$) of (i) follows therefore from \sref{0.22.6.4}, and the equivalence of conditions $a$) and $d$) in (i) follows from \sref{0.22.5.8}; the equivalence of $d$) and $a$) in (ii) follows from \sref{0.22.6.6}; the equivalence of $a$) and $d$) in (ii) follows from the criterion of Zariski \sref{0.22.6.4}, and the last assertion follows from \sref{0.22.6.5}.

Finally, as $\Omega^1_{\mathrm{B}/k_0}$ is a projective $\mathrm{B}$-module, the equivalence of the criterion of Zariski $e$) and $b$) in (ii) follows from the application of \sref{0.19.1.12}, taking into account the fact that $\Omega^1_{\mathrm{B}/\mathrm{A}}$ is a projective $\mathrm{B}$-module \sref{0.20.4.8}, and using \sref{0.22.6.3}, taking into account \sref{0.20.4.8}.
\end{proof}

\begin{corollary}[22.6.8]
\label{0.22.6.8}
\textnormal{(Zariski.)}
Let $\mathrm{C}$ be an algebra of finite type over a field.
Then the set of $\mathfrak{p}\in\Spec(\mathrm{C})$ such that $\mathrm{C}_\mathfrak{p}$ is a regular local ring is open in $\Spec(\mathrm{C})$.
\end{corollary}

\begin{remark}[22.6.9]
\label{0.22.6.9}
\begin{enumerate}
  \item[\rm{(i)}] The statements of \sref{0.22.6.7} are still valid if instead of supposing $k$ separable over $k_0$, one assumes only that it is of \emph{finite radiciel multiplicity}, that is to say \sref{0.19.6.6} that there exists a finite radiciel extension $k_0'$ of $k_0$ such that the residue field of the local Artinian ring $k\otimes_{k_0}k_0'$ is a \emph{separable} extension $k'$ of $k_0'$.
  There exists then a $k_0'$-monomorphism $k'\to k\otimes_{k_0}k_0'$ which, composed with the canonical homomorphism $k\otimes_{k_0}k_0'\to k'$, gives the identity, since $k'$ is a $k_0'$-algebra formally smooth \sref{0.19.6.1}; one concludes that $\mathrm{C}\otimes_{k_0}k_0'$ is equipped with the structure of a $k'$-algebra of finite type, and as $k'$ is separable over $k_0$, one can apply \sref{0.22.6.7} to this $k'$-algebra; as we shall not have occasion to use this generalization, we leave the detail to the reader.
  We do not know whether the results of \sref{0.22.6.7} are valid without \emph{any} hypothesis on the extension $k$ of $k_0$.
  \item[\rm{(ii)}] Under the general hypotheses of \sref{0.22.6.7}, let $(k_\alpha)$ be a decreasing filtering family of subfields of $k$ containing $k_0$ and such that $\bigcap k_\alpha(k^p)=k_0(k^p)$ (where $p$ is the characteristic exponent of $k_0$).
  Using a dimension counting method due to Nagata, one can show that, in the Jacobian criterion \sref{0.22.6.7}, (iii), one may restrict the $k$-derivations $\mathrm{D}_j$ to be $k_\alpha$-derivations for a suitable $\alpha$.
\end{enumerate}
\end{remark}

\subsection{Nagata's Jacobian criterion}
\label{subsection:0.22.7}

\oldpage[0\textsubscript{IV-1}]{209}
\begin{env}[22.7.1]
\label{0.22.7.1}
The Jacobian criterion of Nagata is the analog of the Jacobian criterion of Zariski, but for quotient rings of \emph{formal power series} rings over a field.
We shall give, as Nagata \cite{IV-31}, two versions, presented here as criteria of formal smoothness.
\end{env}

\begin{proposition}[22.7.2]
\label{0.22.7.2}
Let $k$ be a field, $\mathrm{A}=k[[\mathrm{T}_1,\ldots,\mathrm{T}_r]]$ a formal power series ring, equipped with its usual $k$-algebra structure, $\mathfrak{q}$ an ideal of $\mathrm{A}$, $\mathrm{B}=\mathrm{A}/\mathfrak{q}$.
Let $\mathfrak{p}$ be a prime ideal of $\mathrm{A}$ containing $\mathfrak{q}$.
Suppose that there exists a sub-$k$-algebra $\mathrm{C}'$ of $\mathrm{C}=\mathrm{A}/\mathfrak{p}$, isomorphic to a formal power series ring $k[[\mathrm{X}_1,\ldots,\mathrm{X}_s]]$, such that $\mathrm{C}$ is finite over $\mathrm{C}'$ and that the field of fractions $\mathrm{L}$ of $\mathrm{C}$ is a separable extension of $\mathrm{L}'=k((\mathrm{X}_1,\ldots,\mathrm{X}_s))$ (this hypothesis is always satisfied if $k$ is of characteristic $0$).
Then the following conditions are equivalent:
\begin{enumerate}
  \item[a)] $\mathrm{B}_\mathfrak{p}=\mathrm{A}_\mathfrak{p}/\mathfrak{q}\mathrm{A}_\mathfrak{p}$ is a $k$-algebra formally smooth for the $\mathfrak{p}$-preadic topology.
  \item[b)] There exist $k$-derivations $\mathrm{D}_i$ of $\mathrm{A}$ into itself $(1\leq i\leq m)$, and elements $f_i$ of $\mathfrak{q}$ $(1\leq i\leq m)$, such that the images of the $f_i$ in $\mathfrak{q}\mathrm{A}_\mathfrak{p}$ generate this ideal of $\mathrm{A}_\mathfrak{p}$, and that $\det(\mathrm{D}_if_j)\notin\mathfrak{p}$.
  \item[c)] $\mathrm{B}_\mathfrak{p}$ is a regular ring.
\end{enumerate}
\end{proposition}

\begin{proof}
The residue field of $\mathrm{B}_\mathfrak{p}$ is $k(\mathfrak{p})$.
Since $k((\mathrm{X}_1,\ldots,\mathrm{X}_s))$ is separable over $k$ \sref{0.21.9.6.4}, the field $k(\mathfrak{p})=\mathrm{L}$ is separable over $k$ by hypothesis; hence conditions $a$) and $c$) are equivalent by \sref{0.19.6.4}.
Furthermore, $\mathrm{A}_\mathfrak{p}$ is regular \sref{0.17.1.4}, \sref{0.17.3.2}, and therefore formally smooth over $k$ for its $\mathfrak{p}$-preadic topology \sref{0.19.6.4}.
By \sref{0.22.6.2},~(ii), condition $a$) is consequently equivalent to injectivity of
\begin{equation}
\label{0.22.7.2.1}
  j_0:(\mathfrak{q}/\mathfrak{q}^2)\otimes_\mathrm{A}\mathrm{L}
  \longrightarrow \Omega^1_{\mathrm{A}_\mathfrak{p}/k}\otimes_{\mathrm{A}_\mathfrak{p}}\mathrm{L}
  =\widehat{\Omega}^1_{\mathrm{A}/k}\otimes_\mathrm{A}\mathrm{L}.
  \tag{22.7.2.1}
\end{equation}

Consider also
\begin{equation}
\label{0.22.7.2.2}
  j:(\mathfrak{q}/\mathfrak{q}^2)\otimes_\mathrm{A}\mathrm{L}
  \xrightarrow{j_*}\Omega^1_{\mathrm{A}/k}\otimes_\mathrm{A}\mathrm{L}
  \longrightarrow\widehat{\Omega}^1_{\mathrm{A}/k}\otimes_\mathrm{A}\mathrm{L}.
  \tag{22.7.2.2}
\end{equation}
We show first that $b$) is equivalent to injectivity of $j$.
Indeed,
\[
  (\mathfrak{q}/\mathfrak{q}^2)\otimes_\mathrm{A}\mathrm{L}
  =(\mathfrak{q}\mathrm{A}_\mathfrak{p}/\mathfrak{q}^2\mathrm{A}_\mathfrak{p})\otimes_{\mathrm{A}_\mathfrak{p}}\mathrm{L}.
\]
If $F_i=j(f_i'\otimes 1)$, condition $b$) says that the matrix $(\langle F_i,\mathrm{D}_j\otimes 1\rangle)$ is invertible.
Thus the $F_i$, and hence the generators $f_i'\otimes 1$, are linearly independent, so $j$ is injective.
Conversely, if $j$ is injective, choose $f_i$ so that the $f_i'\otimes 1$ form a basis of the source.
Then the $F_i$ are linearly independent.
By \sref{0.21.9.3}, $\widehat{\Omega}^1_{\mathrm{A}/k}$ is free of rank $r$, and the $k$-derivations of $\mathrm{A}$ into itself generate its dual; hence there exist derivations $\mathrm{D}_j$ for which $(\langle F_i,\mathrm{D}_j\otimes 1\rangle)$ is invertible.
This is condition $b$).

These observations show already that $b$) implies $a$).
Conversely, assume $a$), so that $j_0$ is injective, and consider
\begin{equation}
\label{0.22.7.2.3}
\begin{tikzcd}[column sep=large]
  (\mathfrak{q}/\mathfrak{q}^2)\otimes_\mathrm{A}\mathrm{L} \ar[r,"j"] \ar[d,"\alpha"'] & \widehat{\Omega}^1_{\mathrm{A}/k}\otimes_\mathrm{A}\mathrm{L} \ar[d,equal] \\
  (\mathfrak{p}/\mathfrak{p}^2)\otimes_\mathrm{A}\mathrm{L} \ar[r,"i"'] & \widehat{\Omega}^1_{\mathrm{A}/k}\otimes_\mathrm{A}\mathrm{L}.
\end{tikzcd}
\tag{22.7.2.3}
\end{equation}
We shall use the following lemma.
\end{proof}

\begin{lemma}[22.7.2.4]
\label{0.22.7.2.4}
Let $\mathrm{R}$ be a regular local ring, $\mathfrak{m}$ its maximal ideal, $\mathrm{K}=\mathrm{R}/\mathfrak{m}$ its residue field.
If $\mathfrak{n}$ is an ideal of $\mathrm{R}$ such that $\mathrm{R}/\mathfrak{n}$ is regular, then the canonical homomorphism
\[
\label{0.22.7.2.5}
  \alpha:(\mathfrak{n}/\mathfrak{n}^2)\otimes_\mathrm{R}\mathrm{K}\to(\mathfrak{m}/\mathfrak{m}^2)\otimes_\mathrm{R}\mathrm{K}
  \tag{22.7.2.5}
\]
is injective.
\end{lemma}

\begin{proof}
Indeed, one knows \sref{0.17.1.6} that $\mathfrak{n}$ is generated by a subset $(x_i)_{1\leq i\leq l}$ forming part of a regular system of parameters $(x_i)_{1\leq i\leq n}$ of $\mathrm{R}$; the images of the $x_i$ for $1\leq i\leq l$ form a base of $(\mathfrak{n}/\mathfrak{n}^2)\otimes_\mathrm{R}\mathrm{K}$, and their images by $\alpha$ form part of the base of $(\mathfrak{m}/\mathfrak{m}^2)\otimes_\mathrm{R}\mathrm{K}$ formed by the images of the $x_i$ for $1\leq i\leq n$, from which the lemma.
\end{proof}

\begin{proof}[Continuation of the proof of Proposition~{\rm\sref{0.22.7.2}}]
By the lemma and the diagram \sref{0.22.7.2.3}, it remains only to prove that $i$ is injective.
Consider the sequence
\begin{equation}
\label{0.22.7.2.6}
  \mathfrak{p}/\mathfrak{p}^2\xrightarrow{h}
  \widehat{\Omega}^1_{\mathrm{A}/k}\widehat{\otimes}_\mathrm{A}(\mathrm{A}/\mathfrak{p})
  \xrightarrow{g}\widehat{\Omega}^1_{(\mathrm{A}/\mathfrak{p})/k}\to 0.
  \tag{22.7.2.6}
\end{equation}
By \sref{0.20.7.20}, $g$ is surjective and the image of $h$ is dense in $\Ker(g)$ for the $\mathfrak{m}$-adic topology.
Since $\widehat{\Omega}^1_{\mathrm{A}/k}$ is an $\mathrm{A}$-module of finite type, every submodule of
\[
  \widehat{\Omega}^1_{\mathrm{A}/k}\otimes_\mathrm{A}(\mathrm{A}/\mathfrak{p})
  =\widehat{\Omega}^1_{\mathrm{A}/k}\widehat{\otimes}_\mathrm{A}(\mathrm{A}/\mathfrak{p})
\]
is closed for the $\mathfrak{m}$-adic topology \sref[I]{I.7.3.5}; therefore \sref{0.22.7.2.6} is exact.
After tensoring with $\mathrm{L}$, one obtains the exact sequence
\[
  (\mathfrak{p}/\mathfrak{p}^2)\otimes_\mathrm{A}\mathrm{L}\xrightarrow{i}
  \widehat{\Omega}^1_{\mathrm{A}/k}\otimes_\mathrm{A}\mathrm{L}\to
  \widehat{\Omega}^1_{(\mathrm{A}/\mathfrak{p})/k}\otimes_{\mathrm{A}/\mathfrak{p}}\mathrm{L}\to 0.
\]
The separability hypothesis and \sref{0.21.9.5} show that the last module has rank $s$ over $\mathrm{L}$; the middle module has rank $r$ by \sref{0.21.9.3}.
Thus $\rg_\mathrm{L}(\operatorname{Im}i)=r-s$.
Since $\mathrm{A}/\mathfrak{p}$ is finite over a subalgebra isomorphic to $k[[\mathrm{T}_1,\ldots,\mathrm{T}_s]]$, one has $\dim(\mathrm{A}/\mathfrak{p})=s$ by \sref{0.17.1.4} and \sref{0.16.1.5}, and hence
\[
  r-s=\dim(\mathrm{A})-\dim(\mathrm{A}/\mathfrak{p})=\dim(\mathrm{A}_\mathfrak{p})
\]
by \sref{0.16.5.11}.
Finally, $\mathrm{A}_\mathfrak{p}$ is regular, so $\dim(\mathrm{A}_\mathfrak{p})$ is the rank over $\mathrm{L}$ of
$(\mathfrak{p}/\mathfrak{p}^2)\otimes_\mathrm{A}\mathrm{L}$ \sref{0.17.1.1}.
The source and image of $i$ therefore have the same rank, and $i$ is injective.
\end{proof}

\begin{theorem}[22.7.3]
\label{0.22.7.3}
Let $k$ be a field, $\mathrm{A}$ a local Noetherian complete ring with residue field $\mathrm{K}$.
Suppose that:
\begin{enumerate}
  \item[$1^\circ$] $[k:k^p]<+\infty$ (where $p$ is the characteristic exponent of $k$);
  \item[$2^\circ$] $\mathrm{K}$ is a finite extension of a separable extension $\mathrm{K}_0$ of $k$;
  \item[$3^\circ$] $\mathrm{A}$ is equipped with a structure of $\mathrm{K}_0$-algebra formally smooth (for the preadic topology).
\end{enumerate}
Let $\mathfrak{q}$ be an ideal of $\mathrm{A}$, $\mathrm{B}=\mathrm{A}/\mathfrak{q}$, $\mathfrak{p}$ a prime ideal of $\mathrm{A}$ containing $\mathfrak{q}$.
The following conditions are equivalent:
\begin{enumerate}
  \item[a)] $\mathrm{B}_\mathfrak{p}$ is a $k$-algebra formally smooth for the $\mathfrak{p}$-preadic topology.
  \item[b)] There exist $k$-derivations $\mathrm{D}_i$ of $\mathrm{A}$ into itself $(1\leq i\leq m)$, and elements $f_i$ $(1\leq i\leq m)$ of $\mathfrak{q}$, such that the images of the $f_i$ in $\mathfrak{q}\mathrm{A}_\mathfrak{p}$ generate this ideal of $\mathrm{A}_\mathfrak{p}$, and that $\det(\mathrm{D}_if_j)\notin\mathfrak{p}$.
  \item[b')] There exists a sub-extension $k'$ of $\mathrm{K}_0$, containing $\mathrm{K}_0^p$, such that $[\mathrm{K}_0:k']<+\infty$, $k'$-derivations $\mathrm{D}_i$ of $\mathrm{A}$ into itself, and elements $f_i$ of $\mathfrak{q}$ $(1\leq i\leq m)$, such that the images of the $f_i$ in $\mathfrak{q}\mathrm{A}_\mathfrak{p}$ generate this ideal of $\mathrm{A}_\mathfrak{p}$, and that $\det(\mathrm{D}_if_j)\notin\mathfrak{p}$.
\end{enumerate}
\end{theorem}

\begin{proof}
We distinguish two cases, according to whether $p=1$ or $p>1$.

\emph{A)} $k$ \emph{is of characteristic} $0$.
\oldpage[0\textsubscript{IV-1}]{212}
Then $\mathrm{K}$ is a separable extension of $\mathrm{K}_0$, so by \sref{0.19.6.4} that $\mathrm{A}$ is $\mathrm{K}_0$-isomorphic to a formal power series ring $\mathrm{K}[[\mathrm{T}_1,\ldots,\mathrm{T}_r]]$ equipped with its usual $\mathrm{K}$-algebra structure.
But then, taking into account \sref{0.19.8.8}, (ii) applied to $\mathrm{A}/\mathfrak{p}$, replacing $k$ by $\mathrm{K}$, one sees that the general conditions of \sref{0.22.7.2} are satisfied by replacing $k$ by $\mathrm{K}$.
In this case one can immediately apply the conclusions of \sref{0.22.7.2}, and it is immediate that this proves the equivalence of $a$), $b$) (with $k'=\mathrm{K}_0$ in $b'$).

\emph{B)} $k$ \emph{is of characteristic} $p>0$.
As $\mathrm{K}_0$ is separable over $k$, it follows from \sref{0.19.3.5}, (ii) and \sref{0.19.6.1} that $\mathrm{A}$ is a $k$-algebra formally smooth; by virtue of \sref{0.22.5.9}, $\mathrm{A}_\mathfrak{p}$ is also a $k$-algebra formally smooth for the preadic topology.
It follows therefore from \sref{0.22.6.2}, (ii) that the condition $a$) is equivalent to the fact that the homomorphism $j_0$ defined in \sref{0.22.7.2.1} (where $\mathrm{L}=k(\mathfrak{p})$) is injective.
Therefore \sref{0.19.1.12}, $c$) shows already that $b$) implies $a$); and since moreover $b'$) trivially implies $b$), it all amounts to showing that the hypothesis that $j_0$ is injective implies $b'$).

Let us first denote by $k'$ a sub-extension of $\mathrm{K}_0$ such that $[\mathrm{K}_0:k']<+\infty$.
Note that $\widehat{\Omega}^1_{\mathrm{A}/k'}$ is a free $\mathrm{A}$-module of finite type.
By the argument of \sref{0.22.7.2} (for a suitable choice of $k'$), the composite homomorphism
\[
\label{0.22.7.3.1}
  j':(\mathfrak{q}/\mathfrak{q}^2)\otimes_\mathrm{A}\mathrm{L}\xrightarrow{j_0}\Omega^1_{\mathrm{A}/k}\otimes_\mathrm{A}\mathrm{L}\to\widehat{\Omega}^1_{\mathrm{A}/k'}\otimes_\mathrm{A}\mathrm{L}
  \tag{22.7.3.1}
\]
is injective.
Consider again the commutative diagram
\[
\label{0.22.7.3.2}
  \begin{tikzcd}
    (\mathfrak{q}/\mathfrak{q}^2)\otimes_\mathrm{A}\mathrm{L} \ar[r,"j_0"] & \Omega^1_{\mathrm{A}/k}\otimes_\mathrm{A}\mathrm{L} \ar[r] & \widehat{\Omega}^1_{\mathrm{A}/k'}\otimes_\mathrm{A}\mathrm{L}\\
    (\mathfrak{p}/\mathfrak{p}^2)\otimes_\mathrm{A}\mathrm{L} \ar[u,"\alpha"] \ar[r,"i_0"] & \Omega^1_{\mathrm{A}/k}\otimes_\mathrm{A}\mathrm{L} \ar[r] \ar[u,equal] & \widehat{\Omega}^1_{\mathrm{A}/k'}\otimes_\mathrm{A}\mathrm{L} \ar[u,equal]
  \end{tikzcd}
  \tag{22.7.3.2}
\]

Since $j_0$ is injective, we have $\Ker(i_0\circ\alpha)=0$.
But as $\mathrm{A}_\mathfrak{p}$ is regular, and the hypothesis $a$) implies that $\mathrm{B}_\mathfrak{p}=\mathrm{A}_\mathfrak{p}/\mathfrak{q}\mathrm{A}_\mathfrak{p}$ is also regular \sref{0.19.6.5}, the lemma \sref{0.22.7.2.4} shows that $\alpha$ is injective, from which $\alpha^{-1}(\Ker(i_0))=0$.
If $i'$ is the composite homomorphism of the second line of \sref{0.22.7.3.2}, one has similarly
\[
\label{0.22.7.3.3}
  \Ker(i')=\Ker(i).
  \tag{22.7.3.3}
\]

But the same argument as in \sref{0.22.7.2} shows that in \sref{0.22.7.2}, one has an exact sequence
\[
  (\mathfrak{p}/\mathfrak{p}^2)\otimes_\mathrm{A}\mathrm{L}\xrightarrow{i}\widehat{\Omega}^1_{\mathrm{A}/k'}\otimes_\mathrm{A}\mathrm{L}\to\widehat{\Omega}^1_{(\mathrm{A}/\mathfrak{p})/k'}\otimes_{\mathrm{A}/\mathfrak{p}}\mathrm{L}\to 0.
\]

But, the hypothesis of separability made on $\mathrm{L}$, together with \sref{0.21.9.5}, implies, by virtue of \sref{0.21.9.5}, that $\widehat{\Omega}^1_{(\mathrm{A}/\mathfrak{p})/k'}\otimes_{\mathrm{A}/\mathfrak{p}}\mathrm{L}$ is of rank $r$ over $\mathrm{L}$ (where $r-s=\dim(\mathrm{A})-\dim(\mathrm{A}/\mathfrak{p})=\dim(\mathrm{A}_\mathfrak{p})$), from which, using \sref{0.21.9.3},
\[
\label{0.22.7.3.4}
  \rg_\mathrm{L}(\Ker(i'))=\rg_\mathrm{L}\Upsilon_{\mathrm{L}/k}.
  \tag{22.7.3.4}
\]

On the other hand, the residue field of $\mathrm{A}_\mathfrak{p}$ being formally smooth over the prime field \sref{0.20.6.22.1}
\[
  \Upsilon_{\mathrm{L}/k}\xrightarrow{\chi_{\mathrm{A}_\mathfrak{p}/k}}(\mathfrak{p}\mathrm{A}_\mathfrak{p})/(\mathfrak{p}\mathrm{A}_\mathfrak{p})^2\xrightarrow{i}\Omega^1_{\mathrm{A}_\mathfrak{p}/k}\otimes_{\mathrm{A}_\mathfrak{p}}\mathrm{L}
\]
and as $\mathrm{A}_\mathfrak{p}$ is a $k$-algebra formally smooth for the $\mathfrak{p}$-preadic topology, it follows from \sref{0.22.2.2} that $\chi_{\mathrm{A}_\mathfrak{p}/k}$ is an injective homomorphism, hence $\Ker(i)$ is isomorphic to $\Upsilon_{\mathrm{L}/k}$, and taking into account \sref{0.22.7.3.4} and the fact that $\Ker(i)\subset\Ker(i')$, \sref{0.22.7.3.3}, this finishes proving the theorem.
\end{proof}

\oldpage[0\textsubscript{IV-1}]{213}
\begin{remark}[22.7.4]
\label{0.22.7.4}
It was shown in fact (using \sref{0.21.9.6}) that condition $b'$) is even equivalent to the other conditions of \sref{0.22.7.3} when one subjects $k'$ to be one of the fields of a decreasing filtering family $(k_\alpha)$ of subfields of $\mathrm{K}_0$, containing $\mathrm{K}_0^p$, such that $[\mathrm{K}_0:k_\alpha]<+\infty$ for every $\alpha$ and that $\bigcap k_\alpha=\mathrm{K}_0^p$.
\end{remark}

\begin{corollary}[22.7.5]
\label{0.22.7.5}
Under the general hypotheses of \sref{0.22.7.3}, the set of prime ideals $\mathfrak{n}\in\Spec(\mathrm{B})$ such that $\mathrm{B}_\mathfrak{n}$ is a $k$-algebra formally smooth (for the $\mathfrak{n}$-preadic topology) is open in $\Spec(\mathrm{B})$.
\end{corollary}

\begin{proof}
With the notations of \sref{0.22.7.3}, it suffices to see that if $\mathrm{B}_\mathfrak{p}$ is formally smooth, for every $\mathfrak{p}'\in\mathrm{V}(\mathfrak{q})$ sufficiently near to $\mathfrak{p}$ in $\Spec(\mathrm{A})$.
But, if the $k$-derivations $\mathrm{D}_i$ and the elements $f_i$ of $\mathfrak{q}$ satisfy criterion $b$) of \sref{0.22.7.3}, one has also $f=\det(\mathrm{D}_if_j)\notin\mathfrak{p}$ for every $\mathfrak{p}'\notin\mathrm{D}(f)$; on the other hand, there exists $g\notin\mathfrak{p}$ such that the images of $f_i$ in $\mathfrak{q}\mathrm{A}_{\mathfrak{p}'}$ generate this ideal of $\mathrm{A}_{\mathfrak{p}'}$, which finishes the proof of the corollary.
\end{proof}

\begin{corollary}[22.7.6]
\label{0.22.7.6}
\textnormal{(Nagata.)}
Let $\mathrm{B}$ be a local Noetherian complete ring containing a field.
Then the set of $\mathfrak{n}\in\Spec(\mathrm{B})$ such that $\mathrm{B}_\mathfrak{n}$ is regular is open in $\Spec(\mathrm{B})$.
\end{corollary}

\begin{proof}
Indeed, $\mathrm{B}$ is an algebra on a \emph{prime} field $k$, hence \sref{0.19.8.8}, (i) from the proposition $\mathrm{A}/\mathfrak{q}$, where $\mathrm{A}=\mathrm{K}_0[[\mathrm{T}_1,\ldots,\mathrm{T}_n]]$ is a formal power series ring and $\mathrm{K}_0$ an extension of $k$; on the other hand, saying that $\mathrm{B}_\mathfrak{n}$ is regular is equivalent to saying that it is a $k$-algebra formally smooth \sref{0.19.6.4}.
All the conditions of application of \sref{0.22.7.5} are thus realized.
\end{proof}

\begin{remark}[22.7.7]
\label{0.22.7.7}
\begin{enumerate}
  \item[\rm{(i)}] We shall see further on, using \sref{0.22.7.6}, that the conclusion of \sref{0.22.7.6} can be extended to the case of any complete Noetherian local ring (\textbf{IV},~6.12.7).
  \item[\rm{(ii)}] The conclusion of the theorem \sref{0.22.7.3} is no longer necessarily exact when one no longer assumes that $[k:k^p]$ is finite.
  Take for example $k=\mathbf{F}_p(\mathrm{X}_n)_{n\geq 0}$, $\mathrm{A}=k[[\mathrm{T},\mathrm{U}]]$, $\mathfrak{q}$ the principal ideal $\mathrm{A}f$, where $f=\mathrm{U}^p-\sum_{n=0}^\infty\mathrm{X}_n\mathrm{T}^{np}$; $\mathrm{A}$ is a factorial ring in which $f$ is an extremal element, for it is extremal in the polynomial ring $k((\mathrm{T}))[\mathrm{U}]$, hence also in the formal power series ring $k((\mathrm{T}))[[\mathrm{U}]]$ (Bourbaki, \emph{Alg.~comm.}, chap.~VII, \textsection~3), and \emph{a fortiori} in $\mathrm{A}$.
  Take $\mathfrak{p}=\mathfrak{q}$, so that $\mathrm{B}_\mathfrak{p}$ is the field of fractions denoted by $\mathrm{E}$ in \sref{0.21.9.7}, which is separable over $k$ (\emph{loc.~cit.}), hence a $k$-algebra formally smooth.
  But one has, with the notations of \sref{0.22.7.3}, $k=\mathrm{K}_0=\mathrm{K}$, and the $k$-derivations of $\mathrm{A}$ into itself are combinations of $\partial/\partial\mathrm{T}$ and $\partial/\partial\mathrm{U}$; as one has $\partial f/\partial\mathrm{T}=\partial f/\partial\mathrm{U}=0$, the criteria $b$) and $b'$) of \sref{0.22.7.3} are not satisfied.
\end{enumerate}
\end{remark}
